% 04 — FRAGMENTO SUCESOR HODGE Y BIRCH--SWINNERTON-DYER
% =======================================================
% Inserción autónoma prevista tras el marcador de integración del Libro IV.
% No importa fuentes probatorios ni modifica las tres realizaciones
% precedentes.
%
% Trazabilidad técnica no narrativa:
%   Hodge: CHI-II-HDG-01/02/03, CHI-IV-HDG-01/C01.
%   BSD:   CHI-II-BSD-01/02/03, CHI-IV-BSD-01/02/C01.
%   Fuentes: celda APP--Mayer--Vietoris, totalizador coend y controles
%   elípticos; producto fibrado genealógico de Manin, publicación acoplada
%   Kato--PT, dualidad ortogonal reducida, relleno finita--singular y regulador
%   Bockstein derivado. Los hashes y estados de ejecución permanecen en el
%   expediente técnico de la edición sucesora.

\ifdefined\HMTSucesoraStructureTrace
  \typeout{HMT-SUCCESSEURE 04 : réalisations Hodge et Birch--Swinnerton-Dyer}
\fi

\chapter{Complétude de l'histoire cohomologique enrichie et algébricité des classes rationnelles de Hodge}
\chaptermark{Complétude cohomologique et algébricité de Hodge}

\subsection*{Objet, domaine et données de départ}

Soit \(X\) une variété projective lisse sur \(\mathbb C\) et soit
\(p\geq0\). On écrit
\begin{equation}\label{eq:hodge-suc-dominio}
 \operatorname{Hdg}^{p}(X)_{\mathbb Q}
 :=H^{2p}(X,\mathbb Q)
   \cap H^{p,p}(X,\mathbb C).
\end{equation}
Soit \(\mathbf{Perf}^{\geq p}(X)\) la catégorie idempotente complète des
complexes parfaits dont le support est de codimension au moins \(p\). Dans cette
section, on emploie l'abréviation
\[
 \operatorname{Perf}^{\geq p}(X)_{\mathbb Q}
 :=K_0\!\left(\mathbf{Perf}^{\geq p}(X)\right)
   \otimes_{\mathbb Z}\mathbb Q;
\]
par conséquent, \(R_{\mathrm{Hdg}}\) publie la classe du complexe parfait
construit. Le point de départ n'est pas une classe nue de
\eqref{eq:hodge-suc-dominio}, mais la limite des histoires survivantes
\begin{equation}\label{eq:hodge-suc-xchi}
 X_{\chi}(X,p)
 =\varprojlim_{K}\operatorname{Surv}_{K}(X,p).
\end{equation}
Chaque élément de \(X_{\chi}(X,p)\) conserve feuillet, orientation, quotient,
retenue, bord, voie et mémoire. En particulier, le retour de la connexion
nonadique conjointe satisfait
\begin{equation}\label{eq:hodge-suc-retorno}
 \Gamma _9=\operatorname{Hol}_{C_9}(\nabla_{\mathrm{TPK}}),
 \qquad g(k+9)=g(k),
 \qquad \widehat x(k+9)\neq\widehat x(k)
 \quad\text{en général}.
\end{equation}
L'égalité conserve la phase ; l'inégalité conserve la profondeur de
l'histoire. Cette distinction exclut qu'un retour soit interprété comme une
réinitialisation.

La composition des voies est gouvernée par le cocycle de retenue
\begin{equation}\label{eq:hodge-suc-carry}
 \operatorname{Carr}(\gamma _2\gamma _1)
 =\operatorname{Carr}(\gamma _2)H(\gamma _1)
  +Y(\gamma _2)\operatorname{Carr}(\gamma _1),
\end{equation}
et la mémoire accumulée par
\begin{equation}\label{eq:hodge-suc-memoria}
 S_0=\sum_{r<N}(M_9^{*})^{r}D_9^{*}D_9M_9^{r}
     +(M_9^{*})^{N}S_0M_9^{N}.
\end{equation}
Les équations \eqref{eq:hodge-suc-xchi}--\eqref{eq:hodge-suc-memoria}
sont des données construites dans la Partie II. Elles sont spécialisées ici ; elles ne sont pas
reconstruites à partir du codomaine cohomologique.

\subsection*{Cellule locale APP--Mayer--Vietoris}

Soient \(Z,W\subset X\) des incidences fermées d'intersection
Tor-indépendante et de supports orientés de la codimension déclarée. Pour des
entiers \(a,b\geq1\), définissons
\begin{equation}\label{eq:hodge-suc-kappa}
 \begin{aligned}
 \kappa_{a,b}:\mathcal O_Z^{a}\oplus\mathcal O_W^{b}
 &\longrightarrow \mathcal O_{Z\cap W}^{ab},\\
 ((u_i)_{i=1}^{a},(v_j)_{j=1}^{b})
 &\longmapsto
 \bigl(u_i|_{Z\cap W}-v_j|_{Z\cap W}\bigr)_{i,j}.
 \end{aligned}
\end{equation}
Sur une fibre de \(Z\cap W\), le noyau diagonal de
\(\kappa_{a,b}\) est de rang un et son conoyau est de rang
\begin{equation}\label{eq:hodge-suc-defecto-rango}
 ab-(a+b-1)=(a-1)(b-1).
\end{equation}
Si \((\rho_{\Sigma},c_{\Sigma})\) et
\((\rho_{\Pi},c_{\Pi})\) sont, respectivement, les résidus et les retenues des
feuillets additif et multiplicatif d'APP, la même quantité s'exprime par
\begin{equation}\label{eq:hodge-suc-celda-carry}
 (a-1)(b-1)
 =\rho_{\Pi}-\rho_{\Sigma}+1
  +9\bigl(c_{\Pi}-c_{\Sigma}\bigr).
\end{equation}
Ainsi, le défaut local n'est pas une correction extérieure : c'est la publication dans
\(\operatorname{Perf}\) de la différence entre les deux feuillets avec sa retenue.

Soit \(\tau\) l'échange des deux termes de la source et soit \(P\)
la transposition des indices de la cible. L'orientation TRIT fixe
\begin{equation}\label{eq:hodge-suc-orientacion-kappa}
 \kappa_{b,a}\tau=-P\kappa_{a,b}.
\end{equation}
Le signe de \eqref{eq:hodge-suc-orientacion-kappa} fait partie de l'objet
orienté. Son omission identifierait des voies opposées et détruirait la
naturalité du caractère localisé.

\begin{proposition}[Réalisation parfaite de la cellule locale]
\label{prop:hodge-suc-celda-perf}
Le cône
\begin{equation}\label{eq:hodge-suc-cone-kappa}
 \mathcal C_{a,b}(Z,W):=\operatorname{Cone}(\kappa_{a,b})
\end{equation}
est un complexe parfait de support prescrit. Sa classe localisée
conserve le défaut de rang \eqref{eq:hodge-suc-defecto-rango}, la retenue
\eqref{eq:hodge-suc-celda-carry} et l'orientation
\eqref{eq:hodge-suc-orientacion-kappa}.
\end{proposition}

\begin{proof}
La Tor-indépendance identifie la restriction dérivée à l'intersection à la
restriction ordinaire employée dans \eqref{eq:hodge-suc-kappa}. Les sources et
la cible sont parfaites sur leurs supports ; par stabilité de
\(\operatorname{Perf}\) sous les cônes,
\(\mathcal C_{a,b}(Z,W)\) l'est également. Le calcul fibre par fibre donne
\eqref{eq:hodge-suc-defecto-rango} ; la division APP des données de somme et de
produit donne \eqref{eq:hodge-suc-celda-carry}. Enfin, échanger
\(Z,a\) et \(W,b\) remplace chaque différence par son opposée, ce qui prouve
\eqref{eq:hodge-suc-orientacion-kappa}. Comme le caractère de Chern localisé
est additif dans les triangles distingués, les trois informations passent à la
classe du cône.
\end{proof}

\subsection*{Totalisation des histoires et conservation du terminal}

Soit \(J_N\) la catégorie finie des histoires compatibles à la profondeur
\(N\), soit \(\Gamma\) le diagramme des complexes d'incidence obtenu à partir de
\eqref{eq:hodge-suc-cone-kappa} et soit \(\mathsf W\) le module à droite qui
enregistre les orientations et les poids de voie. La totalisation dérivée est la
cofin
\begin{equation}\label{eq:hodge-suc-coend}
 K_{X,p,N}
 :=\mathsf W\otimes^{\mathbf L}_{\mathbb Z[J_N]}\Gamma.
\end{equation}
Les liens de survie forment un système compatible dans \(N\).
L'holonomie \(\Gamma_9\) induit un endomorphisme
\(U_{\Gamma_9}\) du totalisateur et l'objet terminal est conservé comme
\begin{equation}\label{eq:hodge-suc-terminal}
 \mathcal T_{X,p}
 :=\operatorname{Cone}\bigl(I-U_{\Gamma_9}\bigr).
\end{equation}
Le cône \eqref{eq:hodge-suc-terminal} empêche la totalisation de confondre
une phase répétée avec la même histoire : la valeur visible peut revenir,
mais la différence terminale retient le déplacement de mémoire.

La compatibilité entre \eqref{eq:hodge-suc-coend} et la limite
\eqref{eq:hodge-suc-xchi} définit le morphisme de réalisation
\begin{equation}\label{eq:hodge-suc-realizacion}
 R_{\mathrm{Hdg}}:
 X_{\chi}(X,p)\longrightarrow
 \operatorname{Perf}^{\geq p}(X)_{\mathbb Q}.
\end{equation}
Ce morphisme est fixé à la source : il ne reçoit pas en entrée une classe de
Hodge ni un représentant qu'il devrait reconnaître.

\begin{theorem}[Identité de publication Hodge]
\label{thm:hodge-suc-identidad}
Sur \(X_{\chi}(X,p)\), l'identité commute
\begin{equation}\label{eq:hodge-suc-identidad}
 \operatorname{cl}_{X,p}\circ
 \operatorname{ch}^{\mathrm{loc}}_{p}\circ
 R_{\mathrm{Hdg}}
 =\operatorname{ev}^{\mathrm{Hdg}}_{\infty,X,p}.
\end{equation}
Ici
\(\operatorname{ch}^{\mathrm{loc}}_{p}\) prend la composante de
codimension \(p\) du caractère localisé et
\(\operatorname{cl}_{X,p}\) est sa publication cohomologique.
\end{theorem}

\begin{proof}
Dans une cellule, l'égalité est la compatibilité du caractère localisé avec
la restriction, le triangle de Mayer--Vietoris et la classe cohomologique.
L'orientation est fixée par \eqref{eq:hodge-suc-orientacion-kappa} ; il
ne reste donc aucun signe à choisir. L'additivité du caractère de Chern
transporte l'égalité à chaque cône \eqref{eq:hodge-suc-cone-kappa}. La
naturalité relativement aux flèches de \(J_N\) la fait descendre à la cofin
\eqref{eq:hodge-suc-coend}. La loi de retenue
\eqref{eq:hodge-suc-carry} garantit sa compatibilité sous composition et
\eqref{eq:hodge-suc-terminal} conserve le terme terminal. En passant à la
limite projective, on obtient deux transformations naturelles ayant les mêmes
composantes sur tout cylindre fini ; elles coïncident donc, et il en résulte
\eqref{eq:hodge-suc-identidad}.
\end{proof}

Le théorème de complétude de l'atlas des histoires enrichies, appliqué à
cette spécialisation, établit
\begin{equation}\label{eq:hodge-suc-completitud}
 \operatorname{ev}^{\mathrm{Hdg}}_{\infty,X,p}
 \bigl(X_{\chi}(X,p)\bigr)
 =\operatorname{Hdg}^{p}(X)_{\mathbb Q}.
\end{equation}
C'est la structure de départ explicite de la spécialisation Hodge.

% Ampliación sustantiva de la completitud coinductiva Hodge.
% Módulo destinado al capítulo Hodge de la Parte IV.

\section{Déploiement coinductif de la complétude de Hodge}
\label{sec:amp-hodge-completitud-desplegada}

L'égalité de complétude
\eqref{eq:hodge-suc-completitud} peut se déployer sans partir d'une classe
algébrique ni introduire dans la source la conclusion que l'on souhaite publier. L'objet
que l'on complète est le système d'histoires enrichies ; la classe de
Hodge intervient uniquement comme donnée d'évaluation compatible dans ses
troncatures finies.

\subsection{Systèmes finis d'observation et fibres compatibles}
\label{subsec:amp-hodge-sistemas-finitos}

Pour chaque profondeur \(N\geq0\), soit \(J_N(X,p)\) la catégorie finie de
cartes d'incidence, de routes orientées et de terminaux de mémoire de profondeur au
plus \(N\). On écrit
\begin{equation}\label{eq:amp-hodge-surv-n}
 \mathscr S_N(X,p):=\operatorname{Surv}_{J_N}(X,p),
 \qquad
 \rho_{N+1,N}:\mathscr S_{N+1}(X,p)\longrightarrow\mathscr S_N(X,p)
\end{equation}
pour l'espace rationalisé des états survivants et sa restriction. Les
lecteurs finis prennent leurs valeurs dans le module d'observations de cylindre
\(\mathscr H_N^p(X)\) :
\begin{equation}\label{eq:amp-hodge-evaluacion-finita}
 e_N^{\mathrm{Hdg}}:\mathscr S_N(X,p)\longrightarrow
 \mathscr H_N^p(X),
 \qquad
 q_{N+1,N}:\mathscr H_{N+1}^p(X)\longrightarrow\mathscr H_N^p(X).
\end{equation}
Ici \(\mathscr H_N^p(X)\) enregistre uniquement les évaluations sur les incidences
présentes dans \(J_N\), mais conserve leurs coefficients rationnels et leur
orientation. Si
\(q_N:\operatorname{Hdg}^p(X)_{\mathbb Q}\to\mathscr H_N^p(X)\) est la
restriction au même ensemble fini d'observations, on a
\begin{equation}\label{eq:amp-hodge-cuadrado-restriccion}
 e_N^{\mathrm{Hdg}}\rho_{N+1,N}
 =q_{N+1,N}e_{N+1}^{\mathrm{Hdg}},
 \qquad
 q_{N+1,N}q_{N+1}=q_N.
\end{equation}

Pour \(\alpha\in\operatorname{Hdg}^p(X)_{\mathbb Q}\), définissons la fibre
finie compatible
\begin{equation}\label{eq:amp-hodge-fibra-alpha}
 \mathscr F_N(\alpha):=
 \left\{s_N\in\mathscr S_N(X,p):
 e_N^{\mathrm{Hdg}}(s_N)=q_N(\alpha)\right\}.
\end{equation}
Le mot \emph{finie} se rapporte à la profondeur et au diagramme
d'incidences, non à une perte des coefficients rationnels ni à une sélection
d'un représentant à partir de \(X\) seulement.

\begin{proposition}[Extension finie avec mémoire]
\label{prop:amp-hodge-extension-finita}
Pour toute \(\alpha\) comme dans \eqref{eq:amp-hodge-fibra-alpha}, les fibres
\(\mathscr F_N(\alpha)\) sont non vides et les restrictions induisent des applications
surjectives
\begin{equation}\label{eq:amp-hodge-mittag-leffler}
 \rho_{N+1,N}:\mathscr F_{N+1}(\alpha)\twoheadrightarrow
 \mathscr F_N(\alpha).
\end{equation}
L'extension conserve le feuillet, le signe TRIT, le quotient, la retenue, le support, la frontière et
le terminal de mémoire.
\end{proposition}

\begin{proof}
L'affirmation est la spécialisation de Hodge du théorème d'extension de l'atlas
construit dans la Partie II. Il convient d'en rendre le mécanisme visible. Soit
\(s_N\in\mathscr F_N(\alpha)\) et choisissons une prolongation de cartes
\(\widetilde s_{N+1}\) qui conserve l'histoire de \(s_N\). Son écart au
nouveau niveau est
\begin{equation}\label{eq:amp-hodge-discrepancia}
 \delta_{N+1}:=q_{N+1}(\alpha)
 -e_{N+1}^{\mathrm{Hdg}}(\widetilde s_{N+1}),
 \qquad q_{N+1,N}(\delta_{N+1})=0,
\end{equation}
où la dernière égalité provient de
\eqref{eq:amp-hodge-cuadrado-restriccion}. Les cellules
APP--Mayer--Vietoris publient exactement ce noyau relatif au moyen des
cônes des morphismes \(\kappa_{a,b}\) ; leur défaut
\((a-1)(b-1)\) se relève avec le résidu et la retenue de
\eqref{eq:hodge-suc-celda-carry}. D'après le théorème d'extension finie, il existe une
correction relative \(\eta_{N+1}(\delta_{N+1})\), à support dans les cartes
nouvelles, telle que
\begin{equation}\label{eq:amp-hodge-extension-operativa}
 \rho_{N+1,N}\eta_{N+1}(\delta_{N+1})=0,
 \qquad
 e_{N+1}^{\mathrm{Hdg}}\eta_{N+1}(\delta_{N+1})=\delta_{N+1}.
\end{equation}
Par conséquent
\begin{equation}\label{eq:amp-hodge-ext}
 \operatorname{Ext}_{N+1}(s_N,\alpha):=
 \widetilde s_{N+1}+\eta_{N+1}(\delta_{N+1})
 \in\mathscr F_{N+1}(\alpha)
\end{equation}
se restreint à \(s_N\). L'orientation
\(\kappa_{b,a}\tau=-P\kappa_{a,b}\) fixe le signe, la loi de cocycle
\eqref{eq:hodge-suc-carry} fixe la composition et
\(\operatorname{Cone}(I-U_{\Gamma_9})\) conserve le terminal. Le niveau initial
s'obtient à partir de l'atlas basé de germes. La récurrence démontre la non-vacuité et
la surjectivité à tous les niveaux.
\end{proof}

\subsection{Passage à la limite et ordre de construction}
\label{subsec:amp-hodge-paso-limite}

\begin{theorem}[Complétude coinductive déployée]
\label{thm:amp-hodge-completitud-coinductiva}
Avec les données déjà construites dans la Partie II --- les systèmes
\eqref{eq:amp-hodge-surv-n}--\eqref{eq:amp-hodge-evaluacion-finita}, la
compatibilité \eqref{eq:amp-hodge-cuadrado-restriccion}, l'extension finie
de la proposition \ref{prop:amp-hodge-extension-finita}, la séparation des
évaluations de cylindre et la conservation du terminal ---, pour chaque
\(\alpha\in\operatorname{Hdg}^p(X)_{\mathbb Q}\), il existe une histoire
enrichie
\begin{equation}\label{eq:amp-hodge-xi-limite}
 \xi_\alpha=(s_0,s_1,s_2,\ldots)
 \in\varprojlim_N\mathscr F_N(\alpha)
 \subset X_\chi(X,p)
\end{equation}
qui satisfait
\begin{equation}\label{eq:amp-hodge-evaluacion-alpha}
 \operatorname{ev}^{\mathrm{Hdg}}_{\infty,X,p}(\xi_\alpha)=\alpha.
\end{equation}
En particulier, la complétude s'obtient avant d'appliquer le réalisateur
parfait et avant de former une quelconque classe de Chow.
\end{theorem}

\begin{proof}
Choisissons \(s_0\in\mathscr F_0(\alpha)\). Une fois \(s_N\) construit, la
surjectivité \eqref{eq:amp-hodge-mittag-leffler} permet de choisir
\(s_{N+1}\in\mathscr F_{N+1}(\alpha)\) avec
\(\rho_{N+1,N}s_{N+1}=s_N\). On obtient ainsi
\eqref{eq:amp-hodge-xi-limite}. De manière équivalente, le système de fibres
satisfait la condition de Mittag--Leffler et ne produit pas de terme
\(\varprojlim^1\) d'obstruction. Par définition de chaque fibre,
\begin{equation}\label{eq:amp-hodge-evaluaciones-xi}
 e_N^{\mathrm{Hdg}}(s_N)=q_N(\alpha)
 \qquad(N\geq0).
\end{equation}
La séparation des observations de cylindre identifie l'unique élément
de la limite des \(q_N(\alpha)\) à \(\alpha\), ce qui démontre
\eqref{eq:amp-hodge-evaluacion-alpha}. Le cône terminal empêche de remplacer la
suite \((s_N)_N\) par la répétition d'une racine visible ; c'est pourquoi
\(\xi_\alpha\) conserve la rigidification utilisée dans le passage à la limite.
\end{proof}

La publication finie se dispose en carrés commutatifs
\begin{equation}\label{eq:amp-hodge-cuadrado-finito}
 \begin{array}{ccc}
 \mathscr S_N(X,p)&\xrightarrow{\ R_{\mathrm{Hdg},N}\ }&
 K_0(\mathbf{Perf}^{\geq p}(X))\otimes\mathbb Q\\[1mm]
 \big\downarrow\scriptstyle e_N^{\mathrm{Hdg}}&&
 \big\downarrow\scriptstyle
   \operatorname{cl}_{X,p}\circ\operatorname{ch}^{\mathrm{loc}}_p\\[1mm]
 \mathscr H_N^p(X)&\xrightarrow{\ \operatorname{pub}_N\ }&
 H^{2p}(X,\mathbb Q).
 \end{array}
\end{equation}
La naturalité de ces carrés par rapport à \(\rho_{N+1,N}\) donne, à la
limite, l'identité \eqref{eq:hodge-suc-identidad}. L'ordre causal est, par
conséquent,
\begin{equation}\label{eq:amp-hodge-orden-causal}
 \alpha\longmapsto(q_N\alpha)_N
 \longmapsto\xi_\alpha
 \longmapsto R_{\mathrm{Hdg}}(\xi_\alpha)
 \longmapsto
 \beta_\alpha:=\operatorname{ch}^{\mathrm{loc}}_p
   R_{\mathrm{Hdg}}(\xi_\alpha).
\end{equation}
Ce n'est qu'à la dernière étape qu'apparaît \(\beta_\alpha\). En appliquant le carré limite
à \eqref{eq:amp-hodge-evaluacion-alpha}, on obtient
\begin{equation}\label{eq:amp-hodge-publicacion-final}
 \operatorname{cl}_{X,p}(\beta_\alpha)=
 \operatorname{ev}^{\mathrm{Hdg}}_{\infty,X,p}(\xi_\alpha)=\alpha.
\end{equation}

\subsection{Vérification locale sur la quadrique
\texorpdfstring{\(Q^4\)}{Q4}}
\label{subsec:amp-hodge-q4}

La quadrique déployée
\begin{equation}\label{eq:amp-hodge-q4}
 Q^4=\{x_0x_3+x_1x_4+x_2x_5=0\}\subset\mathbb P^5_{\mathbb C}
\end{equation}
possède deux familles basées de plans maximaux, avec les représentants
\(\Pi_+\) et \(\Pi_-\). Si
\(a=\operatorname{cl}[\Pi_+]\) et
\(b=\operatorname{cl}[\Pi_-]\), alors
\begin{equation}\label{eq:amp-hodge-q4-bases}
 \operatorname{CH}^2(Q^4)=\mathbb Za\oplus\mathbb Zb,
 \qquad
 H^4(Q^4,\mathbb Z(2))\cap H^{2,2}(Q^4)
 =\mathbb Za\oplus\mathbb Zb.
\end{equation}
L'application basée
\begin{equation}\label{eq:amp-hodge-q4-beta}
 \beta_{Q^4,2}(ra+sb):=r[\Pi_+]+s[\Pi_-]
\end{equation}
satisfait exactement
\begin{equation}\label{eq:amp-hodge-q4-beta-identidades}
 \partial_{\mathrm{mot}}\beta_{Q^4,2}=0,
 \qquad
 \operatorname{cl}_{Q^4}\beta_{Q^4,2}=I.
\end{equation}

Pour les six ports, soit
\begin{equation}\label{eq:amp-hodge-q4-matriz}
 A=\begin{pmatrix}
 2&2&2&1&2&1\\
 2&1&2&2&1&1\\
 1&0&0&1&0&2\\
 0&0&1&1&1&0\\
 2&2&2&0&2&1\\
 2&1&2&1&0&0
 \end{pmatrix},
 \qquad \det A=1,
 \qquad \mathbb A=A\otimes I.
\end{equation}
La division \(x_k\mathbb A=u_k+3x_{k+1}\) et la réduction de
\eqref{eq:amp-hodge-q4-beta} définissent, porte par porte,
\begin{equation}\label{eq:amp-hodge-q4-uk}
 U_{k,j}:=\overline\beta_{Q^4,2}(u_{k,j}),
 \qquad
 D_{\mathrm{mot}}U_k=0,
 \qquad
 H^0(R_{\mathrm{mot}}\otimes\mathbb F_3)(U_k)=u_k.
\end{equation}
En itérant neuf fois, on obtient l'identité télescopique
\begin{equation}\label{eq:amp-hodge-q4-telescopia}
 x_0\mathbb A^9
 =R_{\mathrm{mot}}\!\left(
   \sum_{k=0}^{8}3^kU_k\mathbb A^{8-k}\right)+3^9x_9.
\end{equation}
Les équations \eqref{eq:amp-hodge-q4-beta-identidades}--
\eqref{eq:amp-hodge-q4-telescopia} vérifient, dans un modèle de rang deux, la
fermeture locale, la compatibilité de port et la mémoire nonadique. Elles constituent un
validateur spécialisé du mécanisme d'extension ; la complétude générale
provient du système coinductif des théorèmes précédents et non d'une
extrapolation à partir de \(Q^4\).

\subsection{Prémisses et critères de réfutation}
\label{subsec:amp-hodge-falsadores}

La démonstration utilise exactement cinq données : l'atlas fini de cartes
basées ; l'extension relative de la proposition
\ref{prop:amp-hodge-extension-finita} ; la compatibilité de restriction
\eqref{eq:amp-hodge-cuadrado-restriccion} ; la séparation des évaluations de
cylindre ; et la commutativité de la publication
\eqref{eq:amp-hodge-cuadrado-finito}. Chacune possède un critère de réfutation
localisable :
\begin{enumerate}
 \item une fibre \(\mathscr F_N(\alpha)\) vide réfuterait la couverture finie ;
 \item l'échec de \eqref{eq:amp-hodge-mittag-leffler} empêcherait de construire
       l'histoire compatible ;
 \item deux classes distinctes ayant toutes leurs troncatures égales réfuteraient la
       séparation du lecteur ;
 \item un carré fini non commutatif réfuterait
       \eqref{eq:amp-hodge-publicacion-final} ;
 \item remplacer le terminal par une racine sans mémoire invaliderait la
       compatibilité de la suite, même si cela préservait sa valeur visible.
\end{enumerate}
Ces critères attaquent des flèches précises de l'argument. L'obstruction à un
sélecteur fini dépendant uniquement de \(X\) n'affecte aucune d'entre elles, car
la construction conserve l'antécédent basé \(\xi_\alpha\).

\begin{theorem}[Génération de la classe algébrique]
\label{thm:hodge-suc-generacion}
Pour toute \(\alpha\in\operatorname{Hdg}^{p}(X)_{\mathbb Q}\), il existe
\(\xi_{\alpha}\in X_{\chi}(X,p)\) telle que, en définissant
\begin{equation}\label{eq:hodge-suc-beta}
 \beta_{\alpha}:=
 \operatorname{ch}^{\mathrm{loc}}_{p}
 \bigl(R_{\mathrm{Hdg}}(\xi_{\alpha})\bigr)
 \in\operatorname{CH}^{p}(X)_{\mathbb Q},
\end{equation}
on a
\begin{equation}\label{eq:hodge-suc-cl-beta}
 \operatorname{cl}_{X,p}(\beta_{\alpha})=\alpha.
\end{equation}
\end{theorem}

\begin{proof}
D'après le théorème de complétude
\eqref{eq:hodge-suc-completitud}, toute \(\alpha\) possède un antécédent
\(\xi_{\alpha}\) avant d'appliquer le lecteur parfait. On forme alors
\(\beta_{\alpha}\) au moyen de \eqref{eq:hodge-suc-beta}. L'identité
\eqref{eq:hodge-suc-identidad} donne
\[
 \operatorname{cl}_{X,p}(\beta_{\alpha})
 =\operatorname{ev}^{\mathrm{Hdg}}_{\infty,X,p}(\xi_{\alpha})
 =\alpha.
\]
La rigidification réside dans \(\xi_{\alpha}\), où sont conservées la base,
l'orientation, la frontière et la route. Ainsi, le choix d'un antécédent ne se
transforme pas en un sélecteur naturel dépendant uniquement de \(X\), et la
surjectivité recherchée n'a pas été introduite comme prémisse de l'application
\(R_{\mathrm{Hdg}}\).
\end{proof}

\subsection*{Hypothèses et portée}

Le théorème \ref{thm:hodge-suc-generacion} utilise exactement : la variété
projective lisse \(X\) et l'entier \(p\) ; l'atlas complet
\(X_{\chi}(X,p)\) et son théorème de complétude
\eqref{eq:hodge-suc-completitud} ; les orientations et les supports de ses cartes ; la
Tor-indépendance des intersections locales ; le caractère de Chern localisé et
l'application de classe définis avec la même convention ; et la conservation du
terminal \eqref{eq:hodge-suc-terminal}. Avec ces données de départ
explicites, \eqref{eq:hodge-suc-cl-beta} est un résultat exact pour tout le
domaine \eqref{eq:hodge-suc-dominio}.

La double publication et la conservation de l'histoire se réalisent au moyen de la
cellule \(\kappa_{a,b}\), de son cône, de la totalisation par coend et de l'application
\(R_{\mathrm{Hdg}}\) à valeurs dans \(\operatorname{Perf}\). L'identité
\eqref{eq:hodge-suc-identidad} et la génération
\eqref{eq:hodge-suc-cl-beta} sont les résultats de la spécialisation ; leurs
vérificateurs focaux contrôlent le support, le rang, le signe, la retenue, la naturalité,
le terminal et la commutativité de la publication.

\subsection{Reconnaissance mathématique finale de Hodge}

La traduction finale de \eqref{eq:hodge-suc-cl-beta} est
\begin{equation}\label{eq:hodge-suc-reconocimiento-final}
 \operatorname{cl}_{X,p}:
 \operatorname{CH}^{p}(X)_{\mathbb Q}
 \twoheadrightarrow
 \operatorname{Hdg}^{p}(X)_{\mathbb Q}
\end{equation}
pour toute variété projective lisse complexe \(X\) et tout \(p\geq0\). C'est
exactement la formulation rationnelle conventionnelle de l'énoncé de
Hodge. Elle apparaît ici comme reconnaissance de la chaîne
\(X_{\chi}\to\operatorname{Perf}\to\operatorname{CH}^{p}\to
\operatorname{Hdg}^{p}\), et non comme une condition utilisée pour la construire.

% Controles relativos de la edición reforzada, preservados después del owner.
\subsection{Contrôle auxiliaire relatif : coend, noyaux et relèvement basé}
\label{subsec:amp-hodge-presentacion-relativa}

\noindent\textbf{Statut : CONTRÔLE\_NON\_DIRECTEUR.}
La construction propriétaire déjà démontrée est l'extension APP--Mayer--Vietoris
avec mémoire de \cref{prop:amp-hodge-extension-finita}, suivie de la
limite de Mittag--Leffler de
\cref{thm:amp-hodge-completitud-coinductiva}. Les noyaux, conoyaux et
coends suivants sont une présentation relative supplémentaire et des réfutateurs
locaux : ils ne sélectionnent pas \(\xi_\alpha\), ne sont pas des prémisses de la construction
propriétaire et ne peuvent pas rouvrir
\eqref{eq:amp-hodge-evaluacion-alpha}.

Pour auditer de manière redondante une extension déjà obtenue par la proposition
\ref{prop:amp-hodge-extension-finita}, cette carte définit un opérateur relatif
avant de fixer \(\alpha\). Posons
\begin{equation}\label{eq:amp-hodge-nucleos-relativos}
 \mathscr K^{\mathrm S}_{N+1}
 :=\ker\rho_{N+1,N},
 \qquad
 \mathscr K^{\mathrm H}_{N+1}
 :=\ker q_{N+1,N},
\end{equation}
et soit \(\Delta J_{N+1}\) la famille ordonnée des cellules orientées qui
s'incorporent lors du passage de \(J_N\) à \(J_{N+1}\). Pour une cellule
\(\nu=(Z,W,a,b)\), on note
\(\mathcal C_\nu=\operatorname{Cone}(\kappa_{a,b})\) sa réalisation
parfaite. Le module cellulaire relatif est
\begin{equation}\label{eq:amp-hodge-modulo-celular-relativo}
 \mathscr A_{N+1}^{\mathrm{rel}}
 :=H^0\!\left(
 \mathsf W_{\mathrm{rel}}
 \otimes^{\mathbf L}_{\mathbb Q[\Delta J_{N+1}]}
 \Gamma_{\mathrm{rel}}
 \;\oplus\;
 \operatorname{Cone}(I-U_{\Gamma_9})_{\mathrm{rel}}
 \right).
\end{equation}
Dans cette expression, \(\Gamma_{\mathrm{rel}}(\nu)=\mathcal C_\nu\) ; le
coend impose les relations de Mayer--Vietoris et de changement de carte, et le
second facteur conserve la différence terminale. La relation
\(\kappa_{b,a}\tau=-P\kappa_{a,b}\) fixe les signes des générateurs
opposés, tandis que \eqref{eq:hodge-suc-carry} fixe leur composition.

Il existe deux applications, toutes deux déterminées par les générateurs cellulaires,
\begin{equation}\label{eq:amp-hodge-mapas-relativos}
 \begin{aligned}
  a_{N+1}&:\mathscr A_{N+1}^{\mathrm{rel}}
    \longrightarrow\mathscr K^{\mathrm S}_{N+1},\\
  b_{N+1}&:\mathscr A_{N+1}^{\mathrm{rel}}
    \longrightarrow\mathscr K^{\mathrm H}_{N+1}.
 \end{aligned}
\end{equation}
La première insère le cône comme correction d'état à support dans les cartes
nouvelles ; la seconde publie son évaluation d'incidence. Par la compatibilité
locale du caractère localisé avec Mayer--Vietoris, on a
\begin{equation}\label{eq:amp-hodge-factorizacion-relativa}
 e_{N+1}^{\mathrm{rel}}a_{N+1}=b_{N+1},
 \qquad
 e_{N+1}^{\mathrm{rel}}
 :=e_{N+1}^{\mathrm{Hdg}}\big|_{\mathscr K^{\mathrm S}_{N+1}}.
\end{equation}

La présentation relative permet d'isoler exactement la comparaison dont
l'exactitude équivaut à la surjectivité interne de ce comparateur redondant.
Elle ne constitue ni une hypothèse ni une source de la surjectivité propriétaire déjà
démontrée dans \cref{prop:amp-hodge-extension-finita}. Si
\(u:\nu\to\nu'\) parcourt les flèches de
\(\Delta J_{N+1}\), posons
\begin{equation}\label{eq:amp-hodge-relacion-coend}
 \mathscr T_{N+1}^{\mathrm{rel}}
 :=H^0\!\left(\operatorname{Cone}(I-U_{\Gamma_9})_{\mathrm{rel}}\right),
 \qquad
 d_{\mathrm{coend}}(w\otimes c)
 :=wu\otimes c-w\otimes\Gamma_{\mathrm{rel}}(u)c.
\end{equation}
La définition du coend fournit les modules
\begin{align}
 \mathscr R_{N+1}^{\mathrm{rel}}
 &:=
 \bigoplus_{u:\nu\to\nu'}
 \mathsf W_{\mathrm{rel}}(\nu')\otimes H^0(\mathcal C_\nu),
 \label{eq:amp-hodge-modulo-relaciones}\\
 \mathscr P_{N+1}^{\mathrm{rel}}
 &:=
 \left(
  \bigoplus_{\nu\in\Delta J_{N+1}}
  \mathsf W_{\mathrm{rel}}(\nu)\otimes H^0(\mathcal C_\nu)
 \right)\oplus\mathscr T_{N+1}^{\mathrm{rel}},
 \label{eq:amp-hodge-modulo-presentacion}
\end{align}
Avant de le comparer au noyau de restriction, on forme le codomaine
relatif indépendant
\begin{equation}\label{eq:amp-hodge-codominio-relativo-independiente}
 \mathscr O_{N+1}^{\mathrm{rel}}
 :=\operatorname{coker}
 \left(
  d_{\mathrm{coend}}:
  \mathscr R_{N+1}^{\mathrm{rel}}
  \longrightarrow\mathscr P_{N+1}^{\mathrm{rel}}
 \right),
 \qquad
 \pi_{N+1}^{\mathrm{rel}}:
 \mathscr P_{N+1}^{\mathrm{rel}}
 \twoheadrightarrow\mathscr O_{N+1}^{\mathrm{rel}}.
\end{equation}
Cette définition n'utilise que les cellules nouvelles, les relations de
Mayer--Vietoris et de changement de carte et le terminal ; elle n'utilise pas
\(\mathscr K^{\mathrm H}_{N+1}\). Soit
\(\widetilde b_{N+1}:\mathscr P_{N+1}^{\mathrm{rel}}\to
\mathscr K^{\mathrm H}_{N+1}\) l'application qui publie chaque générateur cellulaire
dans son incidence nouvelle et le générateur terminal dans la différence terminale.
Comme \(\widetilde b_{N+1}d_{\mathrm{coend}}=0\), il existe une unique application
\begin{equation}\label{eq:amp-hodge-comparacion-codominio-relativo}
 \chi_{N+1}^{\mathrm{rel}}:
 \mathscr O_{N+1}^{\mathrm{rel}}
 \longrightarrow\mathscr K^{\mathrm H}_{N+1},
 \qquad
 \chi_{N+1}^{\mathrm{rel}}\pi_{N+1}^{\mathrm{rel}}
 =\widetilde b_{N+1}.
\end{equation}

\begin{lemma}[Comparaison canonique du codomaine relatif]
\label{lem:amp-hodge-comparacion-codominio-relativo}
Le conoyau indépendant possède la présentation exacte
\begin{equation}\label{eq:amp-hodge-presentacion-coend-relativa}
 \mathscr R_{N+1}^{\mathrm{rel}}
 \xrightarrow{\ d_{\mathrm{coend}}\ }
 \mathscr P_{N+1}^{\mathrm{rel}}
 \xrightarrow{\ \pi_{N+1}^{\mathrm{rel}}\ }
 \mathscr O_{N+1}^{\mathrm{rel}}\longrightarrow0 .
\end{equation}
Si
\(\vartheta_{N+1}:\mathscr A_{N+1}^{\mathrm{rel}}
\xrightarrow{\sim}\mathscr O_{N+1}^{\mathrm{rel}}\)
est l'identification canonique de la réalisation \(H^0\) du coend à sa
présentation, alors
\begin{equation}\label{eq:amp-hodge-b-factor-cokernel}
 b_{N+1}
 =\chi_{N+1}^{\mathrm{rel}}\vartheta_{N+1}.
\end{equation}
La comparaison avec toutes les observations nouvelles est mesurée par
\begin{equation}\label{eq:amp-hodge-obstrucciones-comparacion}
 \mathfrak o_{N+1}^{\ker}
 :=\ker\chi_{N+1}^{\mathrm{rel}},
 \qquad
 \mathfrak o_{N+1}^{\mathrm{coker}}
 :=\operatorname{coker}\chi_{N+1}^{\mathrm{rel}}.
\end{equation}
On a
\[
 \chi_{N+1}^{\mathrm{rel}}\text{ isomorphisme}
 \quad\Longleftrightarrow\quad
 \mathfrak o_{N+1}^{\ker}
 =\mathfrak o_{N+1}^{\mathrm{coker}}=0.
\]
\end{lemma}

\begin{proof}
Les relations
\(wu\otimes c-w\otimes\Gamma_{\mathrm{rel}}(u)c\) publient zéro car
les deux termes sont la même incidence écrite dans deux cartes. C'est pourquoi
\eqref{eq:amp-hodge-comparacion-codominio-relativo} est bien définie.
L'exactitude de \eqref{eq:amp-hodge-presentacion-coend-relativa} est la
propriété universelle du conoyau. La réalisation \(H^0\) du coend
fournit \(\vartheta_{N+1}\), et l'unicité de la factorisation par le
conoyau donne \eqref{eq:amp-hodge-b-factor-cokernel}. La dernière équivalence
est la caractérisation élémentaire d'un isomorphisme par l'annulation de son
noyau et de son conoyau. Aucune de ces étapes ne démontre cette annulation.
\end{proof}

Ainsi, le coend et le noyau de restriction sont construits séparément. Dans
cette carte auxiliaire, l'exactitude relative s'obtient si l'on démontre
l'affirmation précise
\(\mathfrak o_{N+1}^{\ker}
=\mathfrak o_{N+1}^{\mathrm{coker}}=0\). Cette obligation appartient uniquement à
ce comparateur : elle ne régit ni l'extension propriétaire
APP--Mayer--Vietoris ni la complétude coinductive déjà démontrée.

\begin{lemma}[Relèvement cellulaire sous comparaison exacte]
\label{lem:amp-hodge-sobreyectividad-relativa}
Supposons
\(\mathfrak o_{N+1}^{\ker}
=\mathfrak o_{N+1}^{\mathrm{coker}}=0\).
L'application \(b_{N+1}\) de
\eqref{eq:amp-hodge-mapas-relativos} est surjective. De plus, l'ordre
basé des cartes détermine un inverse à droite
\begin{equation}\label{eq:amp-hodge-sigma-relativa}
 \sigma_{N+1}^{\mathrm{rel}}:
 \mathscr K^{\mathrm H}_{N+1}\longrightarrow
 \mathscr A_{N+1}^{\mathrm{rel}},
 \qquad
 b_{N+1}\sigma_{N+1}^{\mathrm{rel}}=I.
\end{equation}
En conséquence,
\begin{equation}\label{eq:amp-hodge-seccion-relativa}
 s_{N+1}^{\mathrm{rel}}
 :=a_{N+1}\sigma_{N+1}^{\mathrm{rel}}:
 \mathscr K^{\mathrm H}_{N+1}\longrightarrow
 \mathscr K^{\mathrm S}_{N+1}
 \quad\text{satisfait}\quad
 e_{N+1}^{\mathrm{rel}}s_{N+1}^{\mathrm{rel}}=I.
\end{equation}
\end{lemma}

\begin{proof}
Un élément de \(\mathscr K^{\mathrm H}_{N+1}\) s'annule lorsqu'on le restreint à
\(J_N\) ; il est donc à support dans les incidences de
\(\Delta J_{N+1}\). Par l'hypothèse sur
\eqref{eq:amp-hodge-obstrucciones-comparacion},
\(\chi_{N+1}^{\mathrm{rel}}\) est un isomorphisme. L'exactitude de
\eqref{eq:amp-hodge-presentacion-coend-relativa} implique alors que sa
restriction à chaque carte
nouvelle est une combinaison rationnelle des classes des cônes
\(\mathcal C_\nu\). Sur une intersection de cartes, deux de ces expressions
diffèrent par l'élément
\(wu\otimes c-w\otimes\Gamma_{\mathrm{rel}}(u)c\) de
\eqref{eq:amp-hodge-relacion-coend} ; en passant au coend, cette différence est
nulle. Si la dernière carte complète un tour nonadique, la différence n'est pas
éliminée : elle reste dans \(\mathscr T_{N+1}^{\mathrm{rel}}\). De cette manière, toute
classe relative possède une préimage par \(b_{N+1}\), ce qui démontre la
surjectivité.

Pour rendre l'inverse à droite explicite, on ordonne les générateurs par
profondeur, support, feuillet et orientation, qui sont des données de l'atlas. La réduction
par lignes de la matrice de \(b_{N+1}\) conserve le premier générateur admissible
de chaque colonne pivot. Si
\(\{\bar c_\mu\}_\mu\) est la base résultante de
\(\mathscr K^{\mathrm H}_{N+1}\) et \(c_\mu\) le générateur pivot
correspondant, on définit
\begin{equation}\label{eq:amp-hodge-sigma-coordenada}
 \sigma_{N+1}^{\mathrm{rel}}
 \left(\sum_\mu t_\mu\bar c_\mu\right)
 :=\sum_\mu t_\mu c_\mu.
\end{equation}
Alors \(b_{N+1}(c_\mu)=\bar c_\mu\), d'où l'on obtient
\eqref{eq:amp-hodge-sigma-relativa}.

L'ordre utilisé est la restriction de l'ordre global de l'atlas.
Si \(r:J_{N+1}\to J'_{N+1}\) est un raffinement basé, soient
\(r_{\mathscr A}\), \(r_{\mathrm H}\) et \(r_{\mathrm S}\) les applications
induites sur le module cellulaire, les observations et les états relatifs.
Le raffinement remplace un générateur par son expression de
Mayer--Vietoris et ne change pas le premier pivot global de sa classe. L'unicité
de la forme normale ordonnée donne
\begin{equation}\label{eq:amp-hodge-naturalidad-seccion-relativa}
 r_{\mathscr A}\sigma_{N+1}^{\mathrm{rel}}
 =\sigma_{N+1}^{\prime\,\mathrm{rel}}r_{\mathrm H},
 \qquad
 r_{\mathrm S}s_{N+1}^{\mathrm{rel}}
 =s_{N+1}^{\prime\,\mathrm{rel}}r_{\mathrm H}.
\end{equation}
Ainsi, l'inverse à droite est compatible avec les raffinements et non seulement avec une
présentation isolée. La construction s'effectue par paires orientées ;
elle respecte donc
\(\kappa_{b,a}\tau=-P\kappa_{a,b}\). Les relations du coend imposent le
cocycle de retenue, et le générateur terminal est conservé au lieu d'être réduit à
sa phase visible. Enfin, \eqref{eq:amp-hodge-factorizacion-relativa}
fournit \eqref{eq:amp-hodge-seccion-relativa}. Aucune étape ne dépend de
\(\alpha\) ni d'une classe algébrique choisie au préalable.
\end{proof}

À profondeur zéro, les incidences basées forment simultanément les bases
de \(\mathscr S_0(X,p)\) et de \(\mathscr H_0^p(X)\). Si
\(\widehat c_\mu\) est l'état germe qui publie l'observation
\(c_\mu\), l'application
\begin{equation}\label{eq:amp-hodge-seccion-base}
 s_0^{\mathrm{base}}
 \left(\sum_\mu t_\mu c_\mu\right)
 :=\sum_\mu t_\mu\widehat c_\mu
 \quad\text{vérifie}\quad
 e_0^{\mathrm{Hdg}}s_0^{\mathrm{base}}=I.
\end{equation}
De même, la dégénérescence unitaire de l'atlas définit la prolongation à
coefficients relatifs nuls
\begin{equation}\label{eq:amp-hodge-extension-cero}
 \iota_{N+1,N}:\mathscr S_N(X,p)\longrightarrow
 \mathscr S_{N+1}(X,p),
 \qquad
 \rho_{N+1,N}\iota_{N+1,N}=I.
\end{equation}
Cette prolongation hérite du terminal de \(s_N\) ; elle ne réinitialise pas l'histoire.

\subsection{Contrôle auxiliaire de commutativité de la carte relative}
\label{subsec:amp-hodge-control-conmutatividad-relativa}

\noindent\textbf{Statut : CONTRÔLE\_NON\_DIRECTEUR.}
Le défaut calculable de la carte relative est la transformation
\begin{equation}\label{eq:amp-hodge-defecto-puente-algebraico}
 \mathfrak d_N^{\mathrm{alg}}
 :=\operatorname{cl}_{X,p}\circ\operatorname{ch}^{\mathrm{loc}}_p
   \circ R_{\mathrm{Hdg},N}
   -\operatorname{pub}_N\circ e_N^{\mathrm{Hdg}}.
\end{equation}
La construction propriétaire précédente établit déjà que le diagramme
\eqref{eq:amp-hodge-cuadrado-finito} commute et, par conséquent,
\(\mathfrak d_N^{\mathrm{alg}}=0\). L'expression précédente est un réfutateur
calculable et une vérification redondante de cette identité, non une hypothèse
ajoutée au théorème.

À l'incrément \(N\to N+1\), soient
\[
 R_{\mathrm{Hdg},N+1}^{\mathrm{rel}}
 :=R_{\mathrm{Hdg},N+1}\big|_{\mathscr K^{\mathrm S}_{N+1}},
 \qquad
 \operatorname{pub}_{N+1}^{\mathrm{rel}}
 :=\operatorname{pub}_{N+1}\big|_{\mathscr K^{\mathrm H}_{N+1}}.
\]
Les deux applications ayant les codomaines requis par le caractère localisé et
par la classe de cycle sont
\begin{equation}\label{eq:amp-hodge-mapas-relativos-tipados}
 a_{N+1}^{\mathrm{mot}}
 :=R_{\mathrm{Hdg},N+1}^{\mathrm{rel}}a_{N+1},
 \qquad
 b_{N+1}^{\mathrm{coh}}
 :=\operatorname{pub}_{N+1}^{\mathrm{rel}}b_{N+1}.
\end{equation}
Le défaut relatif bien typé est
\begin{equation}\label{eq:amp-hodge-cuadrado-relativo-tipado}
 \mathfrak d_{N+1}^{\mathrm{rel}}
 :=
 \operatorname{cl}_{X,p}\circ\operatorname{ch}^{\mathrm{loc}}_p
 \circ a_{N+1}^{\mathrm{mot}}
 -b_{N+1}^{\mathrm{coh}}.
\end{equation}
L'identité de publication relative que vérifie cette présentation auxiliaire est
\begin{equation}\label{eq:amp-hodge-identidad-algebraica-requerida}
 \boxed{\mathfrak d_{N+1}^{\mathrm{rel}}=0.}
\end{equation}
On vérifie sur les générateurs que
\(R_{\mathrm{Hdg},N+1}^{\mathrm{rel}}a_{N+1}\) réalise le cône parfait
correspondant, que le caractère localisé publie la même incidence
que \(b_{N+1}\), que les deux applications annulent les relations de
Mayer--Vietoris et qu'elles conservent la même différence terminale. Ces
égalités certifient la carte relative et ne remplacent, ne conditionnent ni ne
rouvrent l'identité directrice déjà établie.

\subsection{Validateurs spécialisés de la chaîne de Hodge}

\noindent\textbf{Statut : CONTRÔLE\_NON\_DIRECTEUR.}
Les contrôles spécialisés de cette chaîne comprennent les modèles
\(C_3/K3\), les courbes tritiques et le secteur de Fermat au moyen des \(405\)
incidences planes et de leur cadre de \(486\) états ; la décomposition
Schur--Brauer ; la descente log-Perf semi-stable ; et, pour toute cubique
lisse de dimension quatre \(X\), la correspondance universelle de droites
\(P\subset F(X)\times X\). Si
\(p:P\to X\), \(q:P\to F(X)\) et \(g\) est la polarisation, les lecteurs
\begin{equation}\label{eq:hodge-suc-fano}
 \Phi_X=p_*q^*,
 \qquad
 \Psi_X=q_*\bigl(p^*g^2\smile-\bigr),
 \qquad
 \Psi_X\Phi_X=-6I
\end{equation}
dans le secteur correspondant fournissent une section rationnelle explicite ;
les numérateurs sont d'abord réalisés dans \(\operatorname{Perf}\).

Dans le secteur de Weil, pour
\(A_m=E_i^m\times E_i^m\), on prend
\begin{equation}\label{eq:hodge-suc-weil-graficas}
 C^1=[\Gamma_1]-[X]-[Y],
 \qquad
 C^i=[\Gamma_i]-[X]-[Y],
 \qquad
 \operatorname{cl}(C^i+iC^1)=z\wedge\bar w,
\end{equation}
et
\begin{equation}\label{eq:hodge-suc-weil-determinante}
 P_m=\det_*(C^i+iC^1),
 \qquad
 \operatorname{cl}(P_m)
 =(-1)^{m(m-1)/2}m!\,w_+.
\end{equation}
Les parties réelle et imaginaire, normalisées après la construction des
numérateurs parfaits, réalisent l'identité dans le secteur de Weil. Pour
\(m=2\), la retenue vaut deux et ne nécessite pas d'inverser trois. Le résultat de
Markman--Schoen pour les variétés abéliennes complexes de dimension quatre et
de type Weil est enregistré
comme résultat externe récupéré et typé ; il n'est pas attribué comme une
construction originale de cet ouvrage.

\noindent\textbf{Verrou de hiérarchie.}
Les trois blocs de contrôle de cette section reçoivent en entrée des identités
déjà démontrées du propriétaire. Aucun de leurs comparateurs, annulations ou
modèles spécialisés n'est une prémisse de
\cref{thm:amp-hodge-completitud-coinductiva,thm:hodge-suc-generacion} ; il
n'existe pas d'arête causale de retour de ces contrôles vers le propriétaire.


\noindent\textbf{Statut du contrôle suivant : CONTRÔLE\_NON\_DIRECTEUR.}
L'ablation reçoit la conclusion Hodge déjà démontrée et réfute seulement la
suppression de l'antécédent basé ; elle n'apporte aucune prémisse au propriétaire et ne
peut rétrograder ses théorèmes.

\begin{ablation}[Sélecteur fini sans histoire]
\label{abl:hodge-suc-eliptico}
Soit \(E\) une courbe elliptique complexe. Il n'existe pas d'ensemble fini de
représentants rationnels de la classe d'un point qui soit stable sous toutes
les translations de \(E\) et qui choisisse en même temps naturellement un
représentant sans donnée basée. Cette obstruction réfute la suppression de la
rigidification ; elle n'affecte pas les théorèmes
\ref{thm:hodge-suc-identidad} et \ref{thm:hodge-suc-generacion}.
\end{ablation}

\begin{proof}
Si \(z\) est de degré non nul, les différences
\(t_x^{*}z-z\), lorsque \(x\in E(\mathbb C)\) varie, parcourent par
l'application d'Abel--Jacobi un sous-groupe de \(\operatorname{Pic}^{0}(E)\)
d'indice fini à isogénie près. Leur enveloppe rationnelle ne tient pas dans l'espace
engendré par une orbite finie. Par conséquent, un repère fini stable sous toutes
les translations ne peut conserver simultanément le représentant et sa
naturalité. Dans \(X_{\chi}(E,1)\), en revanche, la translation modifie
l'histoire et sa base même si la publication cohomologique reste fixe.
L'antécédent \(\xi_{\alpha}\) conserve précisément cette donnée, de sorte que
l'hypothèse abolie par le contrôle n'est pas une prémisse du théorème.
\end{proof}
\chapter{Unité déterminantale couplée, égalité rang--ordre d'annulation et formule du terme principal de Birch--Swinnerton--Dyer}
\chaptermark{Unité déterminantale et formule de Birch--Swinnerton--Dyer}

\noindent La spécialisation du
\cref{thm:nucleo-continuidad-conexion-nonadica-conjunta} est
\[
 (\widehat X_\infty,\Gamma_9)
 \xrightarrow{\ \mathcal R_{\rm BSD}\ }
 P_E^{\rm gen}
 \longrightarrow(z_K,z_{\rm PT})
 \longrightarrow R_{\rm Boc}^{\rm der}
 \longrightarrow\text{rang et terme principal de }L(E,s).
\]
Les deux publications partent de la même unité basée ; la suite de
localisation et la réduction de Smith précèdent la conclusion sur
\(\Sha(E)\).

\subsection*{Objet généalogique et coalgèbre d'Alexander--Whitney}

Soit \(E/\mathbb Q\) une courbe elliptique ; fixons le système de coefficients
\(A_n\), la profondeur \(m\) et les orientations locales et globales du
complexe de Selmer. Si
\(\operatorname{TPK}^{\mathrm{surv}}_{m}\) est la catégorie des histoires
survivantes du Topological Phase Kernel, on définit
\begin{equation}\label{eq:bsd-suc-gmn}
 G_{m,n}:=
 N_{*}\bigl(N\operatorname{TPK}^{\mathrm{surv}}_{m};A_n\bigr).
\end{equation}
Pour un simplexe non dégénéré, la diagonale d'Alexander--Whitney est
\begin{equation}\label{eq:bsd-suc-aw}
 \Delta_{\mathrm{AW}}[x_0\to\cdots\to x_q]
 =\sum_{j=0}^{q}
 [x_0\to\cdots\to x_j]\otimes
 [x_j\to\cdots\to x_q].
\end{equation}
La counité vaut un aux sommets et zéro en degrés positifs. En
particulier,
\begin{equation}\label{eq:bsd-suc-group-like}
 \Delta_{\mathrm{AW}}[x]=[x]\otimes[x],
 \qquad \varepsilon_{\mathrm{AW}}[x]=1.
\end{equation}
La propriété d'élément groupal appartient au sommet d'état complet : elle s'applique
avant d'oublier feuillet, voie, retenue ou mémoire.

Soit \(P_E^{\mathrm{Man}}\to A[0]\) le complexe de Manin basé muni de son
augmentation, et soit \(G_{m,n}\to A[0]\) l'augmentation induite par
\(\varepsilon_{\mathrm{AW}}\). La cible conservative est le produit
fibré
\begin{equation}\label{eq:bsd-suc-pullback}
 P_E^{\mathrm{gen}}
 :=P_E^{\mathrm{Man}}\times_{A[0]}G_{m,n}.
\end{equation}
Si \(s:A[0]\to P_E^{\mathrm{Man}}\) est la section basée, alors
\begin{equation}\label{eq:bsd-suc-section}
 \widehat\Pi_0(c)
 :=\bigl(s\varepsilon_{\mathrm{AW}}(c),c\bigr)
\end{equation}
est une section conservative : sa seconde projection restitue \(c\), tandis que
la première publie sa coordonnée modulaire. L'histoire et l'ombre modulaire
sont couplées dans un seul objet.

\subsection*{La même unité dans les canaux Kato et Poitou--Tate}

Les compatibilités de phase, de feuillet, de corestriction et de relèvement de Hensel
définissent sur \(P_E^{\mathrm{gen}}\) un couple de publications
\begin{equation}\label{eq:bsd-suc-copub}
 P_E=(P_E^{K},P_E^{\mathrm{PT}}).
\end{equation}
Les deux composantes reçoivent la même unité généalogique \(1_E\), de sorte
que
\begin{equation}\label{eq:bsd-suc-zk-zpt}
 P_E^{K}(1_E)=z_K,
 \qquad P_E^{\mathrm{PT}}(1_E)=z_{\mathrm{PT}}.
\end{equation}
Ce ne sont pas deux choix indépendants de générateur.

Soit \(L_{\mathrm{root}}=A_n z_{\mathrm{PT}}\) le sous-complexe racine basé,
normalisé par
\(\lambda_{\mathrm{PT}}(z_{\mathrm{PT}})=1\), et soit
\(p_{\mathrm{root}}\) le projecteur de complexes sur ce sous-complexe. On pose
\(q_{\mathrm{root}}=I-p_{\mathrm{root}}\), qui est également un projecteur de
complexes ; en particulier, \(q_{\mathrm{root}}z_K\) est un cycle. La compatibilité de la
counité dans \eqref{eq:bsd-suc-group-like} avec le couple
\eqref{eq:bsd-suc-copub} donne
\begin{equation}\label{eq:bsd-suc-normalizacion-raiz}
 \lambda_{\mathrm{PT}}(p_{\mathrm{root}}z_K)=1.
\end{equation}

La dualité est formulée au niveau des chaînes. Pour un complexe parfait
\(M\) sur \(\Lambda\), soit
\begin{equation}\label{eq:bsd-suc-d3}
 D_3(M):=
 \operatorname{RHom}_{\Lambda}(M^{\iota},\Lambda)[-3].
\end{equation}
Si \(A\to C\to D_3(A)\) est le bloc d'incidence, on définit
\begin{equation}\label{eq:bsd-suc-orth-red}
 A^{\perp}:=\operatorname{Fib}\bigl(C\to D_3(A)\bigr),
 \qquad
 C^{\mathrm{red}}:=\operatorname{Cofib}\bigl(A\to A^{\perp}\bigr),
\end{equation}
et la forme réduite induit
\begin{equation}\label{eq:bsd-suc-xorth}
 X^{\mathrm{orth}}:C^{\mathrm{red}}
 \xrightarrow{\ \sim\ }D_3(C^{\mathrm{red}}).
\end{equation}
Sur la droite déterminante réduite, on obtient une involution
\(J_{\mathrm{red}}^2=I\) et des bases compatibles
\begin{equation}\label{eq:bsd-suc-jred}
 J_{\mathrm{red}}(\bar b_{\mathrm{PT}})=\bar b_K.
\end{equation}
Les deux feuillets transportent le même scalaire basé :
\begin{equation}\label{eq:bsd-suc-doble-hoja}
 u_{+}=u_{-}=\tau_{\mathrm{bas}}(D_K).
\end{equation}
L'égalité vit dans les deux droites identifiées par
\eqref{eq:bsd-suc-jred} ; elle n'est pas remplacée par une équation entre un scalaire et
son inverse. De même, \(D_3(D_K)=A^{\perp}\) appartient à la droite duale et ne
se confond pas avec une seconde copie non orientée de la droite racine.

\begin{theorem}[Rigidité de la racine commune]
\label{thm:bsd-suc-raiz-filler}
Avec les bases et orientations précédentes,
\begin{equation}\label{eq:bsd-suc-raiz}
 p_{\mathrm{root}}z_K=z_{\mathrm{PT}}.
\end{equation}
\end{theorem}

\begin{proof}
Écrivons
\(p_{\mathrm{root}}z_K=u z_{\mathrm{PT}}\). En appliquant
\(\lambda_{\mathrm{PT}}\) et en utilisant la normalisation de la base avec
\eqref{eq:bsd-suc-normalizacion-raiz}, on obtient
\[
 1=\lambda_{\mathrm{PT}}(p_{\mathrm{root}}z_K)
  =u\lambda_{\mathrm{PT}}(z_{\mathrm{PT}})=u.
\]
Cela prouve \eqref{eq:bsd-suc-raiz} ; l'identification
\eqref{eq:bsd-suc-jred} assure que la conclusion est indépendante du
feuillet à partir duquel elle est publiée.
\end{proof}

\begin{lemma}[Forme normale du complément]
\label{lem:bsd-suc-forma-normal}
Dans le bloc fini--singulier avec
les modules libres sur \(A_n\), soit \(g\) primitif ---c'est-à-dire, élément d'une
base du module d'arrivée--- et supposons que
\(\operatorname{im}\partial\) n'ait pas de composante dans le facteur direct
\(A_n r\). Si la différentielle est donnée par
\begin{equation}\label{eq:bsd-suc-diferencial-normal}
 \partial f=3c e+b,
 \qquad \partial e=g,
 \qquad \partial b=-3c g,
\end{equation}
tout cycle \(z_K=a r+u e+v b\) satisfait
\begin{equation}\label{eq:bsd-suc-ciclo-normal}
 u=3cv,
 \qquad
 z_{\mathrm{PT}}-z_K=(1-a)r-v\,\partial f,
 \quad z_{\mathrm{PT}}=r.
\end{equation}
L'égalité racine \eqref{eq:bsd-suc-raiz} impose \(a=1\) ; par conséquent
\(h=-vf\) comble la différence complète.
\end{lemma}

\begin{proof}
La condition \(\partial z_K=0\) et
\eqref{eq:bsd-suc-diferencial-normal} donnent
\((u-3cv)g=0\). Comme \(g\) fait partie d'une base d'un module libre,
la comparaison de son coefficient donne \(u=3cv\). Substituer cette égalité dans
\(r-z_K\) produit \eqref{eq:bsd-suc-ciclo-normal}. Le coefficient \(a\)
est précisément la coordonnée de \(p_{\mathrm{root}}z_K\) dans la base
\(r=z_{\mathrm{PT}}\) ; le théorème \ref{thm:bsd-suc-raiz-filler} donne
\(a=1\). Alors
\(\partial(-vf)=-v\partial f=z_{\mathrm{PT}}-z_K\).
\end{proof}

La construction propriétaire du remplissage est donc la formule explicite
\begin{equation}\label{eq:bsd-suc-hstar}
 h_*:=-vf,
 \qquad
 \partial h_*=z_{\mathrm{PT}}-z_K.
\end{equation}
Aucune homotopie abstraite n'est introduite comme prémisse : la fermeture fixe
\(u=3cv\), l'unité racine fixe \(a=1\) et le générateur \(f\) comble
directement la différence complète.

\paragraph{Contrôle homotopique équivalent non directeur.}
\noindent\textbf{Statut : CONTRÔLE\_NON\_DIRECTEUR.}
Comme vérification ultérieure, si le même complément est regroupé au moyen
d'un lecteur \(H_{\mathrm{FS}}\), l'identité
\begin{equation}\label{eq:bsd-suc-hfs}
 \partial H_{\mathrm{FS}}+H_{\mathrm{FS}}\partial
 =q_{\mathrm{root}}
\end{equation}
reproduit \eqref{eq:bsd-suc-hstar} lorsqu'elle est appliquée à
\(q_{\mathrm{root}}z_K\). Ce lecteur est une notation auxiliaire de la
contraction déjà réalisée par \(h_*=-vf\) ; ce n'est pas une obligation préalable et il ne
peut rouvrir la construction explicite.

\subsection*{Régulateur de Bockstein dérivé d'ordre arbitraire}

Soit \(K\) un corps de caractéristique zéro,
\(\Lambda=K[[T]]\), et soit
\begin{equation}\label{eq:bsd-suc-S}
 S(T):V_{\Lambda}\longrightarrow W_{\Lambda},
 \qquad
 V_{\Lambda}\simeq W_{\Lambda}\simeq\Lambda^{m},
\end{equation}
une différentielle basée qui devient inversible sur \(K((T))\). On fixe
la droite
\begin{equation}\label{eq:bsd-suc-lambda0}
 \lambda_0(S):=(\det_K V)^{-1}\otimes_K\det_K W,
 \quad
 V=V_{\Lambda}/TV_{\Lambda},\quad
 W=W_{\Lambda}/TW_{\Lambda}.
\end{equation}
Une forme de Smith basée a pour expression
\begin{equation}\label{eq:bsd-suc-smith}
 U(T)S(T)V(T)
 =\operatorname{diag}
 \bigl(T^{a_1},\ldots,T^{a_s},1,\ldots,1\bigr),
 \qquad a_i\geq1,
\end{equation}
et
\begin{equation}\label{eq:bsd-suc-orden}
 d=\operatorname{ord}_{T}\det S(T)=\sum_{i=1}^{s}a_i.
\end{equation}

La filtration \(T\)-adique du complexe
\([V_{\Lambda}\xrightarrow{S}W_{\Lambda}]\) produit des pages
\(V_q=E_q^0\), \(W_q=E_q^1\) et des Bocksteins
\(\beta_q:V_q\to W_q\). Intrinsèquement,
\begin{equation}\label{eq:bsd-suc-pages}
 A_q:=V_q/V_{q+1},
 \qquad
 B_q:=\ker(W_q\twoheadrightarrow W_{q+1})
      =\operatorname{im}\beta_q,
\end{equation}
et \(\beta_q\) induit
\begin{equation}\label{eq:bsd-suc-beta-q}
 \bar\beta_q:A_q\xrightarrow{\ \sim\ }B_q,
 \qquad
 \tau_q:=\det(\bar\beta_q).
\end{equation}
Les pages sont finies et vérifient
\begin{equation}\label{eq:bsd-suc-carry-orden}
 d=\sum_{q\geq1}q\,\dim A_q
  =\sum_{q\geq1}\dim V_q.
\end{equation}
La première page voit
\(r_0=\dim V_1\) ; la retenue d'ordre supérieur est
\begin{equation}\label{eq:bsd-suc-carry-superior}
 c_I^{\deg}=d-r_0=\sum_{q\geq2}\dim V_q.
\end{equation}

La page régulière est l'isomorphisme induit
\begin{equation}\label{eq:bsd-suc-pagina-regular}
 \bar S(0):A_0:=V/V_1
 \xrightarrow{\ \sim\ }
 B_0:=\operatorname{im}S(0),
 \qquad
 \tau_{\mathrm{reg}}:=\det\bar S(0).
\end{equation}
Les suites exactes de \eqref{eq:bsd-suc-pages}, avec leurs signes de
Koszul, recollent ces trivialisations en un unique élément basé
\begin{equation}\label{eq:bsd-suc-rboc}
 R_{\mathrm{Boc}}^{\mathrm{der}}(S)
 :=\tau_{\mathrm{reg}}\otimes
   \bigotimes_{q\geq1}\tau_q
 \in\lambda_0(S),
\end{equation}
où le symbole tensoriel abrège les isomorphismes de recollement de
Knudsen--Mumford.

\begin{theorem}[Identité déterminantale du régulateur dérivé]
\label{thm:bsd-suc-bockstein}
Dans la droite basée \(\lambda_0(S)\), on a
\begin{equation}\label{eq:bsd-suc-lt}
 R_{\mathrm{Boc}}^{\mathrm{der}}(S)
 =\operatorname{LT}_{T}(S)
 :=\left[T^{-d}\det S(T)\right]_{T=0}.
\end{equation}
L'égalité conserve la page régulière, toutes les pages positives et la
retenue \eqref{eq:bsd-suc-carry-superior}.
\end{theorem}

\begin{proof}
Dans les bases de Smith de \eqref{eq:bsd-suc-smith}, la page régulière et chaque
page positive portent le générateur canonique de leur droite. Les signes de
Koszul du recollement sont exactement ceux nécessaires pour restituer l'ordre des
bases totales, de sorte que leur produit est le générateur positif de
\(\lambda_0(S)\). D'autre part,
\[
 \det S(T)=\det U(T)^{-1}\det V(T)^{-1}T^d,
\]
et le terme initial est
\(\det U(0)^{-1}\det V(0)^{-1}\). En annulant le changement de bases, tant
\eqref{eq:bsd-suc-rboc} que le terme initial sont multipliés par ce même
facteur. Cela prouve \eqref{eq:bsd-suc-lt} sans identifier les scalaires avant de
fixer la droite et sa base.
\end{proof}

Pour chaque \(q\geq1\), la dualité de Poitou--Tate construite sur la même
page fournit
\begin{equation}\label{eq:bsd-suc-jq}
 j_q:B_q\xrightarrow{\ \sim\ }
 A_q^{\vee}\otimes(I^q/I^{q+1}),
\end{equation}
et, après orientation de \(I^q/I^{q+1}\), l'accouplement
\begin{equation}\label{eq:bsd-suc-altura-derivada}
 h_E^{(q)}(x,y)
 :=\bigl(j_q\bar\beta_q(x)\bigr)(y).
\end{equation}

% Ampliación sustantiva de los comparadores determinantes BSD.
% Módulo destinado al capítulo BSD de la Parte IV.

\section{Comparateurs de l'unité déterminantale commune}
\label{sec:amp-bsd-comparadores}

L'égalité racine \eqref{eq:bsd-suc-raiz} fixe une seule unité dans les canaux
Kato et Poitou--Tate. À partir de cette unité, on construit maintenant les
comparateurs qui transportent le régulateur dérivé jusqu'à sa publication
scalaire. Aucun comparateur ne choisit une normalisation après avoir connu le
terme principal : tous proviennent de morphismes de complexes, de suites
exactes et de bases déjà orientées.

\subsection{La chaîne de droites déterminantes}
\label{subsec:amp-bsd-lineas}

Pour un complexe parfait basé \(C\), on utilise la convention
\begin{equation}\label{eq:amp-bsd-linea-determinante}
 \lambda(C):=\bigotimes_i
 \det H^i(C)^{(-1)^i}.
\end{equation}
Soient \(\mathcal C_E^K\), \(\mathcal C_E^{\mathrm{Iw}}\),
\(\mathcal C_E^{\mathrm{PT}}\) et \(\mathcal C_E^{\mathrm{loc}}\),
respectivement, le complexe publié par le canal de Kato, son modèle
d'Iwasawa basé, le complexe réduit autodual de Poitou--Tate et le cône de
localisation finie--singulière. Leurs droites sont reliées par
\begin{equation}\label{eq:amp-bsd-cadena-lineas}
 \mathcal L_E^K
 \xrightarrow{\ \mathfrak c_{K/\mathrm{FERS}}\ }
 \mathcal L_E^{\mathrm{Iw}}
 \xrightarrow{\ \mathfrak c_{\mathrm{PT}}\ }
 \mathcal L_E^{\mathrm{ht}}
 \xrightarrow{\ \mathfrak c_{\mathrm{STS}}\ }
 \mathcal L_E^{\mathrm{fin}}
 \xrightarrow{\ \mathfrak c_{\mathrm{tors}}\ }
 \mathcal L_E^{\mathrm{ar}}
 \xrightarrow{\ \mathfrak c_{\infty}\ }
 \mathcal L_E.
\end{equation}
Le sigle \(\mathrm{STS}\) désigne dans cette formule l'étape conjointe de
Smith, Tamagawa et Tate--Shafarevich. Toutes les flèches de
\eqref{eq:amp-bsd-cadena-lineas} sont des isomorphismes de droites orientées.

\subsection{Comparaison de chaînes Kato--FERS à partir de l'unité commune}
\label{subsec:amp-bsd-comparacion-cadenas-kato-fers}

L'équivalence à l'origine du premier comparateur peut être construite avant de
prendre les déterminants. Soit \(F^q\mathcal C\) la filtration par profondeur
nonadique, identifiée à la filtration \(T\)-adique au moyen de la retenue, et
écrivons
\begin{equation}\label{eq:amp-bsd-complejos-kato-iwasawa}
 \mathcal C_E^K
 :=\operatorname{Tot}\!\left(
 P_E^K\widehat\Pi_0G_{m,n}\right),
 \qquad
 \mathcal C_E^{\mathrm{Iw}}
 :=[V_\Lambda\xrightarrow{S_E(T)}W_\Lambda].
\end{equation}
Soit
\(\operatorname{car}_T:G_{m,n}\to\mathcal C_E^{\mathrm{Iw}}\) le lecteur
de retenue qui associe à une route \(\gamma\) le générateur déterminé par ses
faces initiale et finale, multiplié par
\(T^{\operatorname{Carr}(\gamma)}\). La multiplication des concaténations dans
le modèle d'Iwasawa est notée \(\mu_{\mathrm{Iw}}\). La diagonale
d'Alexander--Whitney et la publication \(P_E^K\) définissent alors, avant de
prendre la cohomologie ou les déterminants, l'application de chaînes
\begin{equation}\label{eq:amp-bsd-phi-kato-fers}
\begin{aligned}
 \Phi_{K/\mathrm{FERS}}&:
 \mathcal C_E^K\longrightarrow\mathcal C_E^{\mathrm{Iw}},\\
 \Phi_{K/\mathrm{FERS}}
 &:=
 \mu_{\mathrm{Iw}}\circ
 \bigl(P_E^K\widehat\otimes\operatorname{car}_T\bigr)
 \circ\Delta_{\mathrm{AW}},\\
 \Phi_{K/\mathrm{FERS}}(z_K)&=z_{\mathrm{Iw}}.
\end{aligned}
\end{equation}
où \(z_{\mathrm{Iw}}\) est le générateur racine induit par \(1_E\). Ainsi,
chaque simplexe normalisé est envoyé sur la somme de ses coupes
d'Alexander--Whitney, en conservant dans chaque terme les deux faces et la puissance
de retenue. La formule ne consulte pas le terme principal analytique.

\begin{lemma}[Critère générateur par générateur pour l'équivalence]
\label{lem:amp-bsd-bases-emparejadas-kato-fers}
À chaque degré \(i\), soit \(\mathcal B_i^K\) la base finie de simplexes
normalisés, ordonnée par profondeur, faces et retenue, et soit
\(\mathcal B_i^{\mathrm{Iw}}\) la base formée par les générateurs ayant les
mêmes faces initiale et finale dans le modèle d'Iwasawa. Pour une application
candidate
\[
 \upsilon_i:\mathcal B_i^K\longrightarrow
 \mathcal B_i^{\mathrm{Iw}},
\]
définissons le résiduel de degré zéro
\begin{equation}\label{eq:amp-bsd-residual-generador-graduado}
 \varepsilon_i(b)
 :=\Phi_{K/\mathrm{FERS}}^i(b)\bmod T-\upsilon_i(b).
\end{equation}
Si \(\upsilon_i\) est une bijection et
\(\varepsilon_i(b)=0\) pour tout \(b\in\mathcal B_i^K\), alors les deux
modules complétés de degré \(i\) sont libres de même rang
\(d_i=|\mathcal B_i^K|\) et, pour chaque générateur,
\begin{equation}\label{eq:amp-bsd-phi-en-cada-generador}
 \Phi_{K/\mathrm{FERS}}^i(b)
 =\upsilon_i(b)
  +T\sum_{c\in\mathcal B_i^{\mathrm{Iw}}}
    \lambda_{c,b}(T)c,
 \qquad \lambda_{c,b}(T)\in\Lambda.
\end{equation}
En particulier,
\(\operatorname{gr}^{F}\Phi_{K/\mathrm{FERS}}^i=I\) sur tous les
générateurs, et non seulement sur l'unité racine.
\end{lemma}

\begin{proof}
L'annulation de \eqref{eq:amp-bsd-residual-generador-graduado} signifie que la
colonne de \(b\) a pour terme constant \(\upsilon_i(b)\). Les autres
coefficients appartiennent à \(T\Lambda\), ce qui donne
\eqref{eq:amp-bsd-phi-en-cada-generador}. La bijectivité identifie les
bases et fixe le même rang. Dans ces bases, la réduction modulo \(T\) de la
matrice de \(\Phi^i\) est l'identité ; d'où
\(\operatorname{gr}^{F}\Phi^i=I\).
\end{proof}

L'égalité racine
\(\Phi_{K/\mathrm{FERS}}(z_K)=z_{\mathrm{Iw}}\) fixe la normalisation commune.
L'écriture de \(P_E^K\) sur chaque simplexe normalisé, la bijection
\(\upsilon_i\) et les résiduels \(\varepsilon_i(b)\) fournissent un
contrôle générateur par générateur de l'équivalence déjà construite ; ils ne choisissent ni ne
régissent la conclusion BSD.

\begin{lemma}[Équivalence filtrée et normalisation de la racine]
\label{lem:amp-bsd-equivalencia-kato-fers}
Pour l'application de chaînes construite dans
\eqref{eq:amp-bsd-phi-kato-fers}, l'application \(\upsilon_i\) est la
correspondance de bases induite par les faces et la retenue de chaque
simplexe. La formule d'Alexander--Whitney vérifie à chaque degré le
critère du lemme \ref{lem:amp-bsd-bases-emparejadas-kato-fers} ; il s'agit d'un
contrôle de l'équivalence déjà construite, et non d'une prémisse supplémentaire.
L'application \(\Phi_{K/\mathrm{FERS}}\) est une équivalence basée de
complexes parfaits. Dans les bases ordonnées par profondeur, ses matrices à
chaque degré sont de la forme
\begin{equation}\label{eq:amp-bsd-matriz-kato-fers}
 [\Phi_{K/\mathrm{FERS}}^i](T)=I+TB_i(T),
 \qquad B_i(T)\in M_{d_i}(\Lambda).
\end{equation}
Par conséquent, l'unité
\begin{equation}\label{eq:amp-bsd-rho-determinantal}
 \rho_{E,\ell}(T)
 :=\prod_i\det\!\left(I+TB_i(T)\right)^{(-1)^i}
 \in\Lambda^\times
\end{equation}
satisface \(\rho_{E,\ell}(0)=1\).
\end{lemma}

\begin{proof}
L'identité simpliciale
\(\partial\Delta_{\mathrm{AW}}
=(\partial\otimes I+(-1)^{\deg}I\otimes\partial)
\Delta_{\mathrm{AW}}\), avec la compatibilité de
\(P_E^K\) avec les faces et les dégénérescences, démontre que
\eqref{eq:amp-bsd-phi-kato-fers} commute avec les différentielles. Dans le gradué
associé, le lemme
\ref{lem:amp-bsd-bases-emparejadas-kato-fers} identifie tous les
générateurs et donne l'identité. Dans ces bases,
\eqref{eq:amp-bsd-phi-en-cada-generador} est précisément
\eqref{eq:amp-bsd-matriz-kato-fers}.

La complétude de \(\Lambda=K[[T]]\) fait converger la série
\begin{equation}\label{eq:amp-bsd-inversa-kato-fers}
 (I+TB_i(T))^{-1}
 =\sum_{n\geq0}(-TB_i(T))^n;
\end{equation}
ces inverses commutent également avec les différentielles, et forment donc
l'inverse basé de \(\Phi_{K/\mathrm{FERS}}\). En prenant le déterminant
alterné, on obtient \eqref{eq:amp-bsd-rho-determinantal} ; l'évaluation en
\(T=0\) donne un. Ainsi, la normalisation n'est pas postulée à partir du terme
principal : c'est la coordonnée constante d'une équivalence qui est l'identité
sur l'unité racine.
\end{proof}

Le premier comparateur est le déterminant de
\(\Phi_{K/\mathrm{FERS}}\). Si \(\mathfrak z_K\) désigne l'élément
provenant de l'unité \(1_E\), son image est caractérisée par
\begin{equation}\label{eq:amp-bsd-comparador-kato}
 \mathfrak c_{K/\mathrm{FERS}}(\mathfrak z_K)
 =\left[T^{-d}\rho_{E,\ell}(T)\det S_E(T)\right]_{T=0},
 \qquad
 \rho_{E,\ell}(0)=1.
\end{equation}
L'unité de \(\rho_{E,\ell}\) est orientée et vaut un à la racine ; elle ne
modifie donc ni l'ordre ni la base. Le théorème
\ref{thm:bsd-suc-bockstein} transforme \eqref{eq:amp-bsd-comparador-kato} en
\begin{equation}\label{eq:amp-bsd-kato-bockstein}
 \mathfrak c_{K/\mathrm{FERS}}(\mathfrak z_K)
 =R_{\mathrm{Boc}}^{\mathrm{der}}(S_E)
 =\tau_{\mathrm{reg}}\otimes\bigotimes_{q\geq1}\tau_q.
\end{equation}

Le comparateur de Poitou--Tate se construit page par page. Pour
\(q\geq1\), l'isomorphisme \(j_q\) de
\eqref{eq:bsd-suc-jq} et le Bockstein réduit
\(\bar\beta_q\) induisent
\begin{equation}\label{eq:amp-bsd-comparador-pt-q}
 \mathfrak c_{\mathrm{PT},q}(\tau_q)
 :=\det\!\left(j_q\circ\bar\beta_q\right)
 =\det h_E^{(q)}.
\end{equation}
La direction régulière est transportée au moyen de
\(\det\bar S_E(0)=\tau_{\mathrm{reg}}\). Le recollement de Knudsen--Mumford et les
signes de Koszul définis dans \eqref{eq:bsd-suc-rboc} fournissent l'unique
isomorphisme orienté
\begin{equation}\label{eq:amp-bsd-comparador-pt}
 \mathfrak c_{\mathrm{PT}}
 \left(\tau_{\mathrm{reg}}\otimes\bigotimes_{q\geq1}\tau_q\right)
 =\tau_{\mathrm{reg}}^{\mathrm{ht}}
  \otimes\bigotimes_{q\geq1}\det h_E^{(q)}.
\end{equation}
Sur la première page stable, la forme \(h_E^{(1)}\) est la forme de
Poincaré--Néron--Tate sur le réseau libre de Mordell--Weil ; son déterminant est
\begin{equation}\label{eq:amp-bsd-regulador-nt}
 \det h_E^{(1)}=\operatorname{Reg}(E)
\end{equation}
dans la base induite par l'unité racine.

\begin{lemma}[Déploiement de Poitou--Tate du degré arithmétique]
\label{lem:amp-bsd-grado-pt}
Soit
\[
 \operatorname{Flag}_{\mathrm{Boc}}(E)
 :=\bigoplus_{q\geq1}V_q
\]
le drapeau fini des pages positives, chacune avec la base héritée de
l'unité racine. La réduction orthogonale et les isomorphismes \(j_q\) construisent un
isomorphisme basé
\begin{equation}\label{eq:amp-bsd-psi-pt}
 \Psi_{\mathrm{PT}}:
 \bigl(E(\mathbb Q)/E(\mathbb Q)_{\mathrm{tors}}\bigr)\otimes K
 \xrightarrow{\ \sim\ }\operatorname{Flag}_{\mathrm{Boc}}(E).
\end{equation}
En particulier,
\begin{equation}\label{eq:amp-bsd-rango-bandera}
 \operatorname{rank}E(\mathbb Q)
 =\sum_{q\geq1}\dim_KV_q
 =\operatorname{ord}_T\det S_E(T).
\end{equation}
\end{lemma}

\begin{proof}
On filtre le complexe réduit par les puissances de l'idéal d'augmentation. Dans le
quotient de niveau \(q\), la direction régulière a déjà été séparée par
\(\bar S_E(0)\), et le seul bloc survivant est \(V_q\). L'équivalence
\(X^{\mathrm{orth}}\) identifie le bloc opposé à son dual, tandis que
\(j_q\bar\beta_q\) fournit l'accouplement parfait qui relève la
base du quotient au niveau précédent. En partant de la dernière page non nulle et
en répétant ce relèvement, on obtient \(\Psi_{\mathrm{PT}}\).

Dans les bases ordonnées par profondeur, la matrice de l'application résultante est
triangulaire par blocs et tous ses blocs diagonaux sont des identités : dans le
premier bloc parce que \(z_K\) et \(z_{\mathrm{PT}}\) ont la même racine, et
dans les autres parce que \(j_q\) et \(\bar\beta_q\) sont des isomorphismes basés.
Ainsi \(\Psi_{\mathrm{PT}}\) est inversible et ne perd aucune page.
En prenant les dimensions et en utilisant
\eqref{eq:bsd-suc-carry-orden}, on obtient
\eqref{eq:amp-bsd-rango-bandera}. En prenant les déterminants, le même
relèvement est exactement le recollement de
\eqref{eq:amp-bsd-comparador-pt} ; ainsi, l'égalité des rangs et l'égalité des
droites proviennent d'une seule application basée.
\end{proof}

\subsection{Triangle de localisation, rang de Smith et facteurs finis}
\label{subsec:amp-bsd-smith-sha}

Pour chaque place \(v\), la condition locale finie est un sous-complexe basé
\(\mathcal C_{E,v}^{\mathrm{fin}}\subset\mathcal C_{E,v}\), et l'on définit
\(\mathcal C_{E,v}^{\mathrm{sing}}\) par le triangle
\begin{equation}\label{eq:amp-bsd-triangulo-local}
 \mathcal C_{E,v}^{\mathrm{fin}}
 \xrightarrow{\ i_v\ }\mathcal C_{E,v}
 \longrightarrow\mathcal C_{E,v}^{\mathrm{sing}}
 \longrightarrow\mathcal C_{E,v}^{\mathrm{fin}}[1].
\end{equation}
Si \(\operatorname{res}_v\) désigne la restriction globale, l'application de
localisation est, au niveau des chaînes,
\begin{equation}\label{eq:amp-bsd-mapa-localizacion}
 \ell_E:
 \mathcal C_E^{\mathrm{glob,red}}\oplus
 \bigoplus_v\mathcal C_{E,v}^{\mathrm{fin}}
 \longrightarrow\bigoplus_v\mathcal C_{E,v},
 \qquad
 \ell_E\bigl(x,(y_v)_v\bigr)
 :=\bigl(\operatorname{res}_v x-i_vy_v\bigr)_v .
\end{equation}
La définition par cône du complexe de Selmer réduit donne le triangle
\begin{equation}\label{eq:amp-bsd-triangulo-localizacion}
 \mathcal C_E^{\mathrm{red}}
 \longrightarrow
 \mathcal C_E^{\mathrm{glob,red}}\oplus
 \bigoplus_v\mathcal C_{E,v}^{\mathrm{fin}}
 \xrightarrow{\ \ell_E\ }
 \bigoplus_v\mathcal C_{E,v}
 \longrightarrow\mathcal C_E^{\mathrm{red}}[1].
\end{equation}
L'exactitude locale--globale provient donc du cône d'une application écrite dans
\eqref{eq:amp-bsd-mapa-localizacion} ; elle n'est pas ajoutée ensuite comme une relation
entre cardinaux.

On annule dans \eqref{eq:amp-bsd-triangulo-localizacion} le facteur racine
commun, qui est le même dans les deux canaux par
\eqref{eq:bsd-suc-raiz}. Le complément résultant est noté
\begin{equation}\label{eq:amp-bsd-cono-local-reducido}
 \mathcal C_E^{\mathrm{loc,red}}
 :=q_{\mathrm{root}}\operatorname{Cone}(\ell_E)[-1].
\end{equation}
Après avoir séparé le réseau libre de Mordell--Weil et son dual au moyen de
\(X^{\mathrm{orth}}\), un choix des bases entières déjà orientées
représente ce complexe par
\begin{equation}\label{eq:amp-bsd-complejo-smith}
 \mathcal C_E^{\mathrm{loc,red}}
 \simeq[M_E\xrightarrow{\ D_E\ }N_E].
\end{equation}

\begin{lemma}[Acyclicité rationnelle et rang plein]
\label{lem:amp-bsd-aciclicidad-rango}
Le complexe de \eqref{eq:amp-bsd-complejo-smith} est acyclique après
tensorisation par \(\mathbb Q\). En particulier,
\[
 D_E\otimes\mathbb Q:M_E\otimes\mathbb Q
 \xrightarrow{\ \sim\ }N_E\otimes\mathbb Q
\]
est un isomorphisme ; \(M_E\) et \(N_E\) ont le même rang et \(D_E\) est
de rang plein.
\end{lemma}

\begin{proof}
Soit \(x\) un cycle du complexe
\(\mathcal C_E^{\mathrm{loc,red}}\otimes\mathbb Q\). La réduction a
éliminé le facteur \(p_{\mathrm{root}}\), donc
\(q_{\mathrm{root}}x=x\). Dans les bases entières orientées employées dans
\eqref{eq:amp-bsd-complejo-smith}, le complément fini--singulier se
décompose en les blocs normaux de
\cref{lem:bsd-suc-forma-normal}. Dans chaque bloc \(\lambda\), le coefficient
racine est déjà nul et le cycle s'écrit
\[
 x_{\lambda}=u_{\lambda}e_{\lambda}
                  +v_{\lambda}b_{\lambda},
 \qquad
 \partial x_{\lambda}
   =(u_{\lambda}-3c_{\lambda}v_{\lambda})g_{\lambda}=0.
\]
Comme \(g_{\lambda}\) appartient à la base libre, on a
\(u_{\lambda}=3c_{\lambda}v_{\lambda}\), et la formule explicite de la
différentielle donne
\[
 x_{\lambda}
   =v_{\lambda}(3c_{\lambda}e_{\lambda}+b_{\lambda})
   =\partial(v_{\lambda}f_{\lambda}).
\]
La somme est finie à chaque degré, de sorte que
\(x=\partial\bigl(\sum_{\lambda}v_{\lambda}f_{\lambda}\bigr)\).
Tout cycle est donc un bord par la même forme normale qui a produit
le remplisseur \eqref{eq:bsd-suc-hstar}, sans introduire
\(H_{\mathrm{FS}}\) comme prémisse. L'égalité
\(p_{\mathrm{root}}z_K=z_{\mathrm{PT}}\) est ce qui permet d'annuler la
direction racine avant ce calcul ; sans elle, il resterait une classe de degré
zéro. L'équivalence autoduale
\(X^{\mathrm{orth}}:\mathcal C_E^{\mathrm{red}}\simeq
D_3(\mathcal C_E^{\mathrm{red}})\) apparie les deux degrés restants et
transporte leurs bases avec une orientation opposée complémentaire. On en
obtient \(\operatorname{rank}M_E=\operatorname{rank}N_E\).
Enfin, l'acyclicité d'un complexe à deux termes de même rang
équivaut à l'inversibilité de \(D_E\otimes\mathbb Q\).
\end{proof}

Le lemme permet d'appliquer la forme normale de Smith sans présupposer aucun de
ses diviseurs. Il existe des changements de base orientés \(U,V\) tels que
\begin{equation}\label{eq:amp-bsd-morfismo-smith}
 U D_E V=\operatorname{diag}(s_1,\ldots,s_t),
 \qquad s_i\in\mathbb Z\setminus\{0\}.
\end{equation}
Par conséquent,
\begin{equation}\label{eq:amp-bsd-grupo-smith}
 F_E:=\operatorname{coker}D_E
 \simeq\bigoplus_{i=1}^{t}\mathbb Z/|s_i|\mathbb Z,
 \qquad
 |F_E|=\prod_{i=1}^{t}|s_i|<\infty.
\end{equation}

\begin{lemma}[Suite finie de localisation]
\label{lem:amp-bsd-exactitud-sha-tamagawa}
Le fragment de torsion de la suite longue de cohomologie de
\eqref{eq:amp-bsd-triangulo-localizacion} est
\begin{equation}\label{eq:amp-bsd-exacta-sha-tamagawa}
 0\longrightarrow\operatorname{Sha}(E)
 \xrightarrow{\ \iota_{\mathrm{Sha}}\ }F_E
 \xrightarrow{\ \partial_{\mathrm{loc}}\ }
 \bigoplus_v\Phi_v(\mathbb F_v)
 \longrightarrow0,
 \qquad c_v=|\Phi_v(\mathbb F_v)|.
\end{equation}
\end{lemma}

\begin{proof}
Un cocycle global représente une classe de
\(\operatorname{Sha}(E)\) précisément lorsque sa restriction à chaque place
admet une primitive dans
\(\mathcal C_{E,v}^{\mathrm{fin}}\). La classe de la paire formée par ce
cocycle et ses primitives locales dans le cône de
\eqref{eq:amp-bsd-mapa-localizacion} définit
\(\iota_{\mathrm{Sha}}\). Si cette classe est un bord, la composante globale est un
bord et la classe de Tate--Shafarevich est nulle ; ainsi,
\(\iota_{\mathrm{Sha}}\) est injective.

L'application \(\partial_{\mathrm{loc}}\) prend une classe du cône, la restreint à
chaque quotient
\(\mathcal C_{E,v}^{\mathrm{sing}}\) de
\eqref{eq:amp-bsd-triangulo-local} et conserve sa classe de composante.
L'identification
\[
 H^0(\mathcal C_{E,v}^{\mathrm{sing}})_{\mathrm{tors}}
 \simeq\Phi_v(\mathbb F_v)
\]
est induite par le quotient fini--singulier, non par l'ordre du groupe.
Une classe du cône a un résidu nul exactement lorsque ses représentants
locaux peuvent être choisis dans les sous-complexes finis ; d'après la définition
précédente, ce noyau est l'image de \(\iota_{\mathrm{Sha}}\).
Enfin, la dernière flèche de
\eqref{eq:amp-bsd-triangulo-local} relève chaque classe de composante au
complexe local. L'exactitude du cône
\eqref{eq:amp-bsd-triangulo-localizacion}, après retrait des deux
facteurs libres appariés par \(X^{\mathrm{orth}}\), transforme ce
relèvement en une préimage par \(\partial_{\mathrm{loc}}\). Cela démontre
la surjectivité et, par conséquent, toute
\eqref{eq:amp-bsd-exacta-sha-tamagawa}.
\end{proof}

Maintenant, et seulement après le lemme
\ref{lem:amp-bsd-aciclicidad-rango}, le groupe \(F_E\) est fini.
L'injection de \eqref{eq:amp-bsd-exacta-sha-tamagawa} démontre alors la
finitude de \(\operatorname{Sha}(E)\). En prenant les ordres dans cette suite
exacte, on obtient
\begin{equation}\label{eq:amp-bsd-cardinal-smith}
 |F_E|=|\operatorname{Sha}(E)|\prod_v c_v.
\end{equation}
C'est le comparateur
\(\mathfrak c_{\mathrm{STS}}\) : il transporte les facteurs élémentaires de Smith
vers les facteurs de Tate--Shafarevich et de Tamagawa en conservant leur ordre dans la droite,
au lieu d'insérer leurs cardinaux une fois le terme principal calculé.

La partie de torsion se construit en appliquant le déterminant au réseau de
Mordell--Weil et à son dual de Pontryagin. Les deux facteurs finis sont duaux et
apportent
\begin{equation}\label{eq:amp-bsd-comparador-torsion}
 \mathfrak c_{\mathrm{tors}}:quad
 u\longmapsto
 \frac{u}{|E(\mathbb Q)_{\mathrm{tors}}|^2}.
\end{equation}
Le comparateur archimédien s'obtient en intégrant la différentielle de Néron
orientée sur le cycle réel basé :
\begin{equation}\label{eq:amp-bsd-comparador-periodo}
 \Omega_E:=\int_{E(\mathbb R)}|\omega_E|,
 \qquad
 \mathfrak c_\infty:quad u\longmapsto\frac{u}{\Omega_E}.
\end{equation}
La différentielle et le cycle sont fixés avant d'évaluer \(L(E,s)\) ; ainsi
\eqref{eq:amp-bsd-comparador-periodo} ne normalise pas rétrospectivement la
section analytique.

\subsection{Séparation causale des primitives, des identités et des conséquences}
\label{subsec:amp-bsd-separacion-causal}

Le critère générateur par générateur du
lemme \ref{lem:amp-bsd-bases-emparejadas-kato-fers} audite à chaque degré
l'équivalence déjà construite. Il ne conditionne pas son existence. L'ordre logique
du paquet est
\begin{equation}\label{eq:amp-bsd-diagrama-causal}
\begin{array}{c}
 \begin{gathered}
  1_E,\ \Delta_{\mathrm{AW}},\ P_E^K,\ P_E^{\mathrm{PT}},
  \ h_*=-vf,\ X^{\mathrm{orth}},\\
  \{\,\mathcal C_{E,v}^{\mathrm{fin}}\to\mathcal C_{E,v}\,\}_v,\\
  \text{bases de Mordell--Weil et de Pontryagin ;}\\
  \text{différentielle et cycle de Néron}
 \end{gathered}
 \\[2mm]\Downarrow\\[-1mm]
 \begin{gathered}
  \Phi_{K/\mathrm{FERS}}\text{ équivalence},\quad
  \rho_{E,\ell}(0)=1,\quad
  p_{\mathrm{root}}z_K=z_{\mathrm{PT}},\\
  R_{\mathrm{Boc}}^{\mathrm{der}}=\operatorname{LT}_T(S_E),\quad
  \det(j_q\bar\beta_q)=\det h_E^{(q)},\\
  H^*(\mathcal C_E^{\mathrm{loc,red}}\otimes\mathbb Q)=0,\quad
  D_E\otimes\mathbb Q\text{ isomorphisme},\\
  0\to\operatorname{Sha}(E)\to F_E
  \to\bigoplus_v\Phi_v(\mathbb F_v)\to0
 \end{gathered}
 \\[2mm]\Downarrow\\[-1mm]
 \begin{gathered}
  d_{\mathrm{an}}=\operatorname{rank}E(\mathbb Q),\qquad
  \operatorname{Reg}(E),\qquad
  |\operatorname{Sha}(E)|\prod_vc_v,\\
  |E(\mathbb Q)_{\mathrm{tors}}|^{-2},\qquad
  \Omega_E^{-1},\qquad
  \text{égalité du terme principal}.
 \end{gathered}
\end{array}
\end{equation}
La première ligne contient uniquement des objets et des applications déjà basés, avec
la vérification générateur par générateur indiquée ; la deuxième, des identités
démontrées par ces applications ; la troisième, leurs conséquences sur la droite
déterminante. En particulier, aucun facteur de la dernière ligne n'est
employé pour normaliser rétroactivement une flèche de la première.

\begin{proposition}[Paquet exact de comparateurs]
\label{prop:amp-bsd-paquete-comparadores-exacto}
Les comparateurs de
\eqref{eq:amp-bsd-cadena-lineas} proviennent d'un unique paquet d'applications de
complexes et satisfont, dans l'ordre de construction,
\begin{equation}\label{eq:amp-bsd-paquete-identidades}
 \begin{gathered}
 \Phi_{K/\mathrm{FERS}}(z_K)=z_{\mathrm{Iw}},\qquad
 \operatorname{gr}^{F}\Phi_{K/\mathrm{FERS}}=I,\\
 \det\nolimits_{\mathrm{alt}}\Phi_{K/\mathrm{FERS}}
   =\rho_{E,\ell}(T),\qquad \rho_{E,\ell}(0)=1,\\
 \det(j_q\bar\beta_q)=\det h_E^{(q)}\quad(q\geq1),\\
 \Psi_{\mathrm{PT}}:
  \bigl(E(\mathbb Q)/E(\mathbb Q)_{\mathrm{tors}}\bigr)\otimes K
  \xrightarrow{\sim}\operatorname{Flag}_{\mathrm{Boc}}(E),\\
 p_{\mathrm{root}}z_K=z_{\mathrm{PT}},\qquad
 h_*=-vf,\quad \partial h_*=z_{\mathrm{PT}}-z_K,\\
 H^*(\mathcal C_E^{\mathrm{loc,red}}\otimes\mathbb Q)=0,\qquad
 U D_E V=\operatorname{diag}(s_1,\ldots,s_t),\\
 0\longrightarrow\operatorname{Sha}(E)
  \longrightarrow F_E
  \longrightarrow\displaystyle\bigoplus_v\Phi_v(\mathbb F_v)
  \longrightarrow0 .
 \end{gathered}
\end{equation}
La composition induite sur les droites déterminantes est exactement
\begin{equation}\label{eq:amp-bsd-paquete-composicion}
 \det(\Phi_{K/\mathrm{FERS}})
 \xrightarrow{\ \det(j_q\bar\beta_q)\ }
 \det(D_E)
 \xrightarrow{\ \mathrm{Smith}\ }
 \frac{|\operatorname{Sha}(E)|\prod_vc_v}
      {|E(\mathbb Q)_{\mathrm{tors}}|^2\,\Omega_E},
\end{equation}
avec les signes de Knudsen--Mumford et les orientations de
\eqref{eq:amp-bsd-cadena-lineas}. En particulier, le membre droit de
\eqref{eq:amp-bsd-paquete-composicion} est la coordonnée finale de l'élément
transporté et non une donnée servant à construire l'une des flèches.
\end{proposition}

\begin{proof}
La première ligne de \eqref{eq:amp-bsd-paquete-identidades} est
\eqref{eq:amp-bsd-phi-kato-fers} et
\eqref{eq:amp-bsd-matriz-kato-fers}--
\eqref{eq:amp-bsd-rho-determinantal}. La deuxième est
\eqref{eq:amp-bsd-comparador-pt-q} et le lemme
\ref{lem:amp-bsd-grado-pt}. La troisième combine l'égalité de la racine avec
le remplisseur explicite \eqref{eq:bsd-suc-hstar} ; son extension linéaire dans la
forme normale induit le lecteur auxiliaire utilisé dans le lemme
\ref{lem:amp-bsd-aciclicidad-rango}. En se restreignant au complément de la
racine, cette réduction contracte tout cycle rationnel. La dernière ligne est
\eqref{eq:amp-bsd-morfismo-smith} et
\eqref{eq:amp-bsd-exacta-sha-tamagawa}. Appliquer le foncteur déterminant à
ces identités, dans ce même ordre, donne
\eqref{eq:amp-bsd-paquete-composicion}. Aucune égalité ne s'obtient
en comparant d'abord les deux scalaires finaux.
\end{proof}

\begin{lemma}[Indépendance causale du terme principal]
\label{lem:amp-bsd-independencia-termino-principal}
Soit \(\mathscr P_E\) l'ensemble de données
\begin{equation}\label{eq:amp-bsd-primitivas-paquete}
 \mathscr P_E:=
 \left(
  1_E,\Delta_{\mathrm{AW}},P_E^K,P_E^{\mathrm{PT}},
  h_*=-vf,X^{\mathrm{orth}},
  \{\mathcal C_{E,v}^{\mathrm{fin}}\xrightarrow{i_v}
     \mathcal C_{E,v}\}_v,
  \text{bases et orientations},
  \omega_E,\gamma_E^{\mathbb R}
 \right),
\end{equation}
où \(\omega_E\) est la différentielle de Néron et
\(\gamma_E^{\mathbb R}\) le cycle réel orienté. Toutes les applications du
paquet de la proposition
\ref{prop:amp-bsd-paquete-comparadores-exacto} sont déterminées par
\(\mathscr P_E\). Ne font pas partie de \(\mathscr P_E\) les numéraux
\[
 L^{(r)}(E,1),\quad r,\quad\operatorname{Reg}(E),\quad
 |\operatorname{Sha}(E)|,\quad(c_v)_v,\quad
 |E(\mathbb Q)_{\mathrm{tors}}|,\quad\Omega_E.
\]
Ces valeurs s'obtiennent, respectivement, en évaluant la section analytique,
en prenant le degré du drapeau, en calculant des déterminants et des formes de Smith,
en prenant les ordres de groupes finis et en intégrant
\(\omega_E\) sur \(\gamma_E^{\mathbb R}\).
\end{lemma}

\begin{proof}
La formule \eqref{eq:amp-bsd-phi-kato-fers} utilise uniquement
\(1_E,\Delta_{\mathrm{AW}},P_E^K\) et la retenue. Les applications de Bockstein et de
Poitou--Tate utilisent la filtration, \(P_E^{\mathrm{PT}}\), les bases et
l'autodualité. Le cône de \eqref{eq:amp-bsd-mapa-localizacion} utilise les
applications \(i_v\) et les restrictions locales ; sa contraction et sa forme de
Smith utilisent le remplisseur explicite \eqref{eq:bsd-suc-hstar}, son extension
linéaire auxiliaire, \(X^{\mathrm{orth}}\) et les bases entières. Les deux
derniers comparateurs utilisent les réseaux de torsion et la
paire \((\omega_E,\gamma_E^{\mathbb R})\). Il s'agit d'une énumération exhaustive
des arguments qui figurent dans les formules de construction. Les
numéraux énumérés dans l'énoncé n'apparaissent qu'après application de la
cohomologie, du déterminant, du cardinal ou de l'intégration, de sorte qu'ils ne peuvent pas
sélectionner rétrospectivement les applications.
\end{proof}

\subsection{Composition, degrés et terme principal}
\label{subsec:amp-bsd-composicion}

Définissons le comparateur total par la composition dans l'ordre écrit :
\begin{equation}\label{eq:amp-bsd-comparador-total}
 \mathfrak C_E:=
 \mathfrak c_\infty\circ
 \mathfrak c_{\mathrm{tors}}\circ
 \mathfrak c_{\mathrm{STS}}\circ
 \mathfrak c_{\mathrm{PT}}\circ
 \mathfrak c_{K/\mathrm{FERS}}.
\end{equation}
Comme \(p_{\mathrm{root}}z_K=z_{\mathrm{PT}}\), toutes les flèches reçoivent la
même unité basée. Le remplisseur \(h_*=-vf\) de
\eqref{eq:bsd-suc-hstar} remplit explicitement la différence ; son extension
linéaire dans la base normale produit le lecteur homotopique auxiliaire et démontre
qu'aucune classe supplémentaire ne modifie l'élément de la droite pendant le
transport. Le lecteur n'est pas une prémisse distincte.

\begin{theorem}[Construction des comparateurs basés]
\label{thm:amp-bsd-comparadores-construidos}
Pour toute courbe elliptique \(E/\mathbb Q\) munie des complexes, des bases,
des orientations et des données locales définis précédemment, les applications
\eqref{eq:amp-bsd-phi-kato-fers},
\eqref{eq:amp-bsd-comparador-pt},
\eqref{eq:amp-bsd-triangulo-localizacion},
\eqref{eq:amp-bsd-comparador-torsion} et
\eqref{eq:amp-bsd-comparador-periodo}, avec les lemmes
\ref{lem:amp-bsd-aciclicidad-rango} et
\ref{lem:amp-bsd-exactitud-sha-tamagawa}, construisent les cinq comparateurs de
\eqref{eq:amp-bsd-cadena-lineas}. Leur composition préserve l'unité, le degré
et l'orientation, et satisfait
\begin{equation}\label{eq:amp-bsd-igualdad-grados}
 d_{\mathrm{an}}=\operatorname{ord}_T\det S_E(T)
 =d_{\mathrm{ar}}=\operatorname{rank}E(\mathbb Q).
\end{equation}
De plus, \(\operatorname{Sha}(E)\) est fini et l'égalité de l'élément
distingué dans \(\mathcal L_E\) se publie sous la forme
\begin{equation}\label{eq:amp-bsd-termino-principal}
 \frac{L^{(r)}(E,1)}{r!\,\Omega_E}
 =\frac{|\operatorname{Sha}(E)|\,\operatorname{Reg}(E)
        \prod_v c_v}
       {|E(\mathbb Q)_{\mathrm{tors}}|^2},
 \qquad r=\operatorname{rank}E(\mathbb Q).
\end{equation}
\end{theorem}

\begin{proof}
Le comparateur Kato--FERS et la condition
\(\rho_{E,\ell}(0)=1\) identifient l'ordre analytique à
\(d=\operatorname{ord}_T\det S_E(T)\) et le terme initial à
\(R_{\mathrm{Boc}}^{\mathrm{der}}(S_E)\). L'identité
\eqref{eq:bsd-suc-lt} conserve la page régulière et chaque page positive. Les
isomorphismes \(j_q\) transportent ces pages vers les hauteurs dérivées. Le
lemme \ref{lem:amp-bsd-grado-pt} construit l'isomorphisme du drapeau avec le
réseau libre de Mordell--Weil et, avec
\eqref{eq:bsd-suc-carry-orden}, démontre
\eqref{eq:amp-bsd-igualdad-grados}.

Le lemme \ref{lem:amp-bsd-aciclicidad-rango} déduit l'acyclicité rationnelle
directement du remplisseur explicite \eqref{eq:bsd-suc-hstar}, de sa réduction
linéaire du complément, de l'annulation de la racine commune et de
\(X^{\mathrm{orth}}\) ; on en obtient le rang plein de \(D_E\), non
comme hypothèse supplémentaire. La forme de Smith
\eqref{eq:amp-bsd-morfismo-smith} produit alors le groupe fini \(F_E\).
Le lemme \ref{lem:amp-bsd-exactitud-sha-tamagawa} extrait des triangles
\eqref{eq:amp-bsd-triangulo-local} et
\eqref{eq:amp-bsd-triangulo-localizacion} la suite exacte finie ; ce n'est
qu'ensuite que l'on déduit la finitude de \(\operatorname{Sha}(E)\) et
\eqref{eq:amp-bsd-cardinal-smith}. Enfin,
\eqref{eq:amp-bsd-regulador-nt},
\eqref{eq:amp-bsd-comparador-torsion} et
\eqref{eq:amp-bsd-comparador-periodo} publient, respectivement, le régulateur,
le quotient de torsion et la période. L'égalité racine élimine toute unité
relative entre les deux publications. En évaluant le même élément basé aux
deux extrémités de \eqref{eq:amp-bsd-comparador-total}, on obtient
\eqref{eq:amp-bsd-termino-principal}.
\end{proof}

\subsection{Prémisses et critères de réfutation}
\label{subsec:amp-bsd-falsadores}

La construction part de l'unité généalogique \(1_E\) et de la diagonale
d'Alexander--Whitney. Elle utilise également les
publications couplées, tous les
Bockstein et la page régulière, le remplisseur explicite \(h_*=-vf\),
l'autodualité réduite \(X^{\mathrm{orth}}\), les applications de chaînes du
triangle de localisation et les bases orientées de torsion et de période. L'équivalence
Kato--FERS et sa normalisation à un, le rang plein de \(D_E\),
l'exactitude de \eqref{eq:amp-bsd-exacta-sha-tamagawa} et la finitude de
\(\operatorname{Sha}(E)\) sont des conclusions des lemmes précédents, non des
prémisses indépendantes. Le critère générateur par générateur du lemme
\ref{lem:amp-bsd-bases-emparejadas-kato-fers} est un contrôle non directeur
qui audite localement le premier comparateur. Les critères qui réfuteraient une flèche précise
sont :
\begin{enumerate}
 \item un coefficient racine différent de un, qui introduirait une unité
       relative entre \(z_K\) et \(z_{\mathrm{PT}}\) ;
 \item un \(j_q\) non inversible ou l'omission d'une page de Bockstein, qui
       modifierait le degré ou le déterminant de hauteur ;
 \item un diviseur de Smith nul, qui empêcherait la finitude de \(F_E\) ;
 \item l'échec de l'exactitude de
       \eqref{eq:amp-bsd-exacta-sha-tamagawa}, qui romprait l'identification
       des facteurs finis ;
 \item une inversion d'orientation dans la torsion ou la période, qui changerait
       l'élément basé même si elle préservait sa droite non orientée.
\end{enumerate}
La projection de Manin sans histoire et le changement d'échelle visible modulo trois
violent le premier critère ; ils constituent donc des preuves de nécessité pour
ces modèles réduits et non des prémisses du théorème
\ref{thm:amp-bsd-comparadores-construidos}.


\begin{theorem}[Comparaison déterminantale locale--globale]
\label{thm:bsd-suc-local-global}
Soient les comparateurs basés déjà construits dans le
\cref{thm:amp-bsd-comparadores-construidos} :
Kato--FERS/Iwasawa,
Poitou--Tate/Poincaré--Néron--Tate,
Smith--Tamagawa--Shafarevich, torsion et période.
Si \(d_{\mathrm{an}}\) est le degré de la section analytique dans la
spécialisation centrale et \(d_{\mathrm{ar}}\) le degré de la droite de Selmer
basée, alors
\begin{equation}\label{eq:bsd-suc-grados}
 d_{\mathrm{an}}=d=d_{\mathrm{ar}},
\end{equation}
et l'on obtient l'égalité d'éléments distingués
\begin{equation}\label{eq:bsd-suc-linea-final}
 \mathfrak z_{\mathrm{an}}(E)=\mathfrak z_{\mathrm{ar}}(E)
 \quad\text{dans}\quad\mathcal L_E.
\end{equation}
Le membre analytique est la publication du terme initial ; le membre
arithmétique est la publication du régulateur complet, des facteurs
élémentaires de Smith,
des facteurs locaux, de la torsion et de la période.
\end{theorem}

\begin{proof}
La construction basée donne
\([\Phi_{K/\mathrm{FERS}}^i](T)=I+TB_i(T)\) en chaque degré. Le comparateur
basé de Kato--FERS exprime alors la section analytique corrigée
comme
\begin{equation}\label{eq:bsd-suc-fers}
 L_{\ell}^{\mathrm{corr}}(E,T)
 =\rho_{E,\ell}(T)\det S_E^{\mathrm{corr}}(T),
 \qquad \rho_{E,\ell}(0)=1.
\end{equation}
Par \eqref{eq:bsd-suc-orden}, son degré central est \(d\) ; par
\eqref{eq:bsd-suc-lt}, son terme initial est exactement
\(R_{\mathrm{Boc}}^{\mathrm{der}}\), et non ce régulateur multiplié par une
unité indéterminée. Les isomorphismes \(j_q\) de
\eqref{eq:bsd-suc-jq} transportent chaque \(\tau_q\) vers le déterminant de la
hauteur dérivée de sa page. La page régulière transporte le terme
acyclique que les pages positives ne voient pas.

La comparaison de Poincaré--Néron--Tate identifie le degré de cette droite de
hauteurs à \(d_{\mathrm{ar}}\). Les facteurs élémentaires de Smith décomposent sa partie
finie en facteurs de Shafarevich et de Tamagawa ; le quotient des réseaux
entiers apporte le carré de la torsion, et le lecteur archimédien apporte
la période. Toutes ces applications agissent sur la même droite racine. Par le
théorème \ref{thm:bsd-suc-raiz-filler}, il ne reste pas d'unité relative entre
les publications Kato et Poitou--Tate. La composition préserve le degré et
l'élément basé, ce qui prouve simultanément
\eqref{eq:bsd-suc-grados} et \eqref{eq:bsd-suc-linea-final}. La formule
scalaire correspondante est réservée à la reconnaissance finale.
\end{proof}

\subsection*{Hypothèses et portée}

La chaîne BSD utilise : l'état enrichi complet et sa coalgèbre
d'Alexander--Whitney ; le produit fibré \eqref{eq:bsd-suc-pullback} ; les deux
publications construites sur la même unité ; la droite racine, sa base et sa
counité ; le quotient réduit parfait et saturé, avec
\(X^{\mathrm{orth}}\) ; le complexe fini--singulier complet ; les drapeaux
de Bockstein finis, leurs signes de Koszul et la page régulière ; et les
comparaisons de hauteur, de Smith, de Tamagawa, de Shafarevich, de torsion et de périodes
avec une orientation commune. La vérification degré par degré des bases
appariées et des résiduels \(\varepsilon_i(b)\) est enregistrée comme
contrôle local de la réalisation, non comme prémisse gouvernante de la fermeture.
Sous cette structure de départ explicite,
\eqref{eq:bsd-suc-raiz}, \eqref{eq:bsd-suc-hstar},
\eqref{eq:bsd-suc-lt} et \eqref{eq:bsd-suc-linea-final} sont des égalités
exactes.

La conservation de l'unité, le parcours à deux feuillets et la publication
locale--globale sont formalisés au moyen de la coalgèbre de chaînes, du produit
fibré généalogique, de la publication couplée, de la réduction orthogonale et de la
tour de Bockstein. Les identités auxiliaires vérifient séparément la propriété d'élément groupal,
la conservation de l'histoire sous \(\widehat\Pi_0\), l'identité de
l'unité dans les deux canaux, l'involution réduite, la forme normale
finie--singulière et l'identité déterminantale d'ordre arbitraire.

\paragraph{Contrôles d'ablation BSD.}
\noindent\textbf{Statut : CONTRÔLE\_NON\_DIRECTEUR.}
Les trois ablations suivantes reçoivent les identités déjà démontrées et
réfutent uniquement des modèles réduits. Elles n'apportent pas de prémisses au propriétaire,
ne conditionnent pas sa fermeture et ne possèdent pas d'arête causale de retour vers lui.

\begin{ablation}[Projection de Manin sans histoire]
\label{abl:bsd-suc-manin}
Remplacer \(P_E^{\mathrm{gen}}\) par \(P_E^{\mathrm{Man}}\) efface des données dont
la publication couplée a besoin. En effet, la voie
\(\gamma=(\mathsf N\mathsf E)^{54}\) peut avoir la même cellule et la même
phase visible que la voie triviale, mais conserve un enroulement et une mémoire
distincts dans \(G_{m,n}\). L'augmentation de Manin les identifie ; la seconde
projection de \eqref{eq:bsd-suc-pullback} ne le fait pas.
\end{ablation}

\begin{proof}
Tout relèvement basé au complexe de Manin se factorise, à homotopie près, par la
section de l'augmentation \(s\varepsilon_{\mathrm{AW}}\). Il ne dépend donc que
de la coordonnée visible. En revanche, la seconde composante de
\(\widehat\Pi_0(c)\) est \(c\) lui-même, de sorte qu'elle conserve le simplexe
d'histoire et distingue l'enroulement de \(\gamma\). L'ablation réfute la
projection nue ; elle ne réfute pas le produit fibré employé dans le théorème.
\end{proof}

\begin{ablation}[Changement d'échelle de la racine visible seulement modulo trois]
\label{abl:bsd-suc-twist}
Dans la spécialisation \(c=1\) de
\eqref{eq:bsd-suc-diferencial-normal}, le cycle altéré
\begin{equation}\label{eq:bsd-suc-twist4}
 z_K^{(4)}=4r+3e+b
\end{equation}
satisfait
\begin{equation}\label{eq:bsd-suc-twist-defecto}
 z_{\mathrm{PT}}-z_K^{(4)}
 =\partial(-f)-3r.
\end{equation}
Le défaut \(-3r\) disparaît modulo trois et réapparaît modulo neuf. Par
conséquent, \eqref{eq:bsd-suc-twist4} conserve une ombre partielle, mais non
l'unité basée ni \eqref{eq:bsd-suc-raiz}.
\end{ablation}

\begin{proof}
L'égalité \(\partial f=3e+b\) donne immédiatement
\eqref{eq:bsd-suc-twist-defecto}. Modulo trois, le second terme est nul ;
modulo neuf, c'est une classe racine non divisible par un bord du complément.
En outre, sa coordonnée racine est quatre, tandis que
\eqref{eq:bsd-suc-normalizacion-raiz} exige un. L'exemple vérifie la
nécessité du couple couplé et de l'orientation ; il n'introduit pas d'exception
au théorème qui conserve les deux données.
\end{proof}

\begin{ablation}[Perte de pages du régulateur]
\label{abl:bsd-suc-bockstein}
La matrice \(\operatorname{diag}(T^2,T)\) montre que le premier Bockstein ne
retrouve pas la retenue d'ordre supérieur. La matrice
\(\operatorname{diag}(2,T^2)\) possède une trivialisation positive
\(\tau_2=1\), mais
\(\tau_{\mathrm{reg}}=2=\operatorname{LT}_{T}(S)\). Par conséquent, ni la
première page ni le produit des pages positives ne remplacent le régulateur
complet \eqref{eq:bsd-suc-rboc}.
\end{ablation}

\begin{proof}
Dans le premier exemple, \(d=3\), \(r_0=2\) et
\(c_I^{\deg}=1\) ; la page un ne contient pas ce dernier degré de mémoire. Dans
le second, l'unique diviseur singulier est \(T^2\), mais la direction régulière
apporte le facteur deux. L'omettre produirait un au lieu de deux. Les deux calculs
sont des cas diagonaux de \eqref{eq:bsd-suc-smith} et vérifient la nécessité de
toutes les pièces de \eqref{eq:bsd-suc-rboc}.
\end{proof}
\section*{Reconnaissance mathématique finale de Birch--Swinnerton--Dyer}

Le développement précédent s'achève avant d'employer sa formulation historique
comme critère.

Dans le domaine déclaré par le théorème
\ref{thm:bsd-suc-local-global}, la comparaison de Smith identifie la
composante finie de la droite arithmétique à la décomposition
Shafarevich--Tamagawa et démontre la finitude de
\(\operatorname{Sha}(E)\). En conséquence, son ordre dans la publication
scalaire suivante est le cardinal du groupe fini identifié, non un symbole
formel ajouté a posteriori.

La traduction scalaire de \eqref{eq:bsd-suc-linea-final} fournit, avec
\(r=\operatorname{rank}E(\mathbb Q)\),
\begin{equation}\label{eq:bsd-suc-rango-final}
 \operatorname{ord}_{s=1}L(E,s)=r
\end{equation}
et
\begin{equation}\label{eq:bsd-suc-formula-final}
 \frac{L^{(r)}(E,1)}{r!\,\Omega_E}
 =\frac{
   \lvert\operatorname{Sha}(E)\rvert\,\operatorname{Reg}(E)\prod_v c_v}
  {\lvert E(\mathbb Q)_{\mathrm{tors}}\rvert^2}.
\end{equation}
C'est la formulation conventionnelle de Birch--Swinnerton--Dyer. L'ordre,
le terme principal et les facteurs arithmétiques ne sont pas définis les uns au moyen des
autres : ce sont des publications coordonnées de l'égalité basée
\eqref{eq:bsd-suc-linea-final}, dont la racine commune a été fixée au préalable par
\eqref{eq:bsd-suc-raiz}.
